\documentclass[10pt,a4paper]{article}
\usepackage{amsmath,amssymb,amsthm,graphicx}
\usepackage{hyperref}
\usepackage{fontspec}
\title{Plane local testing of multiplicity codes via algebraic plane characterization}
\author{Dan Karliner, Roie Salama, Amnon Ta-Shma}
\date{doi:10.4230/LIPIcs.CCC.2022.14; CC BY 4.0}
\begin{document}
\maketitle
\noindent\textit{Editorial scope.} Counted original Theorem64 (restating Theorem9), the quantitative plane tester for multiplicity codes. Complete Theorem60 plane-characterization argument and all local supporting algebraic and sampling proofs are bundled, not separately counted. Full author RM-to-MRM transfer proof is necessary context. Field-characteristic, degree and multiplicity constraints remain exactly as published. No author theorem, notation or proof is mathematically corrected.\par\medskip
\noindent\textit{Allowed standard prerequisites (editorial).}\begin{itemize}
\item Classical Grobner basis theory, Combinatorial Nullstellensatz with multiplicity [3], Lucas theorem, and imported Reed-Muller local tester [12]. Exact author imported statements, definitions and their source-local explanations are retained. All local representing-polynomial, coefficient-isolation, degree and plane-sampling proofs retained; no recursive reconstruction of cited background results.
\end{itemize}\medskip
% Faithful official CCC2022 article14 transcription: necessary definitions and complete Section5 transfer theorem.
\section*{2 Preliminaries (author context excerpts)}
\subsection*{2.1 Notation}
We denote vectors by bold letters. For $\boldsymbol X=X_1,\ldots,X_m$ we denote by $\mathbb F[\boldsymbol X]$ the set of multivariate polynomials in the variables $X_1,\ldots,X_m$. We denote by $\mathbb F[\boldsymbol X]^{\leq d}$ the set of polynomials of individual degree at most $d$, and by $\mathbb F[\boldsymbol X]^{\mathrm{loc}\leq d}$ the set of polynomials of individual degree at most $d$ (i.e degree in each variable). Given a vector $\boldsymbol I=(i_1,\ldots,i_m)\in\mathbb N^m$, we use the notation
\[\boldsymbol X^{\boldsymbol I}\stackrel{\mathrm{def}}=\prod_{k=1}^mX_k^{i_k}.\]
\subsection*{2.2 Reed Muller and Multiplicity codes}
\paragraph{Definition 10.} Let $d,m$ be non-negative integers, and $q$ a prime power. The $(m,d,q)$ Reed-Muller code is defined as the set of evaluation vectors of $m$-variate polynomials of degree $\leq d$ over $\mathbb F_q^m$, namely,
\[\mathrm{RM}(q,m,d)=\{(f(\alpha))_{\alpha\in\mathbb F_q^m}\mid f\in\mathbb F_q[\boldsymbol X]^{\leq d}\}.\]
We will make use of the following lemma:
\paragraph{Lemma 11 ([10], Lemma 3.2).} Let $\delta_{q,m,d}^{\mathrm{RM}}$ be the relative distance of the code $\mathrm{RM}(q,m,d)$. Then $\delta_{q,m,d}^{\mathrm{RM}}\geq q^{-d/(q-1)}$.
\paragraph{Definition 12 (Hasse derivative).} For a multivariate $P(\boldsymbol X)\in\mathbb F[\boldsymbol X]$ where $\boldsymbol X=(X_1,\ldots,X_m)$ for some $m\in\mathbb N$, and a non-negative vector $\boldsymbol I\in\mathbb N^m$, the $\boldsymbol I$-th Hasse derivative of $P$, denoted $P^{(\boldsymbol I)}(\boldsymbol X)$, is the coefficient of $\boldsymbol Z^{\boldsymbol I}$ in the polynomial $P(\boldsymbol X,\boldsymbol Z)=P(\boldsymbol X+\boldsymbol Z)$. Thus
\[P(\boldsymbol X+\boldsymbol Z)=\sum_{\boldsymbol I}P^{(\boldsymbol I)}(\boldsymbol X)\cdot\boldsymbol Z^{\boldsymbol i}.\]
Hasse derivatives are linear. I.e, for all $P,Q\in\mathbb F[\boldsymbol X]$ and $\lambda\in\mathbb F$, $(\lambda P)^{(\boldsymbol I)}(\boldsymbol X)=\lambda P(\boldsymbol X)$ and $P^{(\boldsymbol I)}(\boldsymbol X)+Q^{(\boldsymbol I)}(\boldsymbol X)=(P+Q)^{(\boldsymbol I)}(\boldsymbol X)$. The product rule shows $(PQ)^{(\boldsymbol I)}(\boldsymbol X)=\sum_{\boldsymbol I_0+\boldsymbol I_1=\boldsymbol I}P^{(\boldsymbol I_0)}(\boldsymbol X)\cdot Q^{\boldsymbol I_1}(\boldsymbol X)$.
\paragraph{Definition 13 (Weight).} If $\boldsymbol I=(i_1,\ldots,i_m)\in\mathbb N^m$ then $w(\boldsymbol I)=\sum_{j=1}^mi_j$.
\paragraph{Definition 14 (Multiplicity).} For $P(\boldsymbol X)\in\mathbb F[\boldsymbol X]$ and $\boldsymbol a\in\mathbb F^m$, the multiplicity of $P$ at $\boldsymbol a$, denoted $\operatorname{mult}(P,\boldsymbol a)$, is the largest integer $s$ such that for every non-negative vector $\boldsymbol I$ with $w(\boldsymbol I)<s$ we have $P^{(\boldsymbol I)}(\boldsymbol a)=0$. If $s$ may be taken arbitrarily large, we set $\operatorname{mult}(P,\boldsymbol a)=\infty$.
Note that by definition $\operatorname{mult}(P,\boldsymbol a)\geq0$ for every $\boldsymbol a$. One important property about multiplicities is a generalization of the Schwartz-Zippel lemma for multivariate polynomials:
\paragraph{Lemma 15 ([6]).} Let $P\in\mathbb F[\boldsymbol X]$ be a non-zero polynomial of total degree at most $d$. Then for any finite $A\subseteq\mathbb F$,
\[\sum_{\boldsymbol a\in A^m}\operatorname{mult}(P,\boldsymbol a)\leq d\cdot|A|^{m-1}.\]
\paragraph{Definition 16 (Multiplicity code).} Let $m,d\geq s$ be non-negative integers, and let $q$ be a prime power. Let
\[\Sigma_{m,s}=\mathbb F_q^{\{\boldsymbol I:w(\boldsymbol I)<s\}}\simeq\mathbb F_q^{\binom{m+s-1}{m}}.\]
For $P(\boldsymbol X)\in\mathbb F_q[X_1,\ldots,X_m]$ we define the order $s$ evaluation of $P$ at $\boldsymbol a$, denoted $P^{(<s)}(\boldsymbol a)$, to be the vector $(P^{(\boldsymbol I)}(\boldsymbol a))_{\boldsymbol I:w(\boldsymbol I)<s}\in\Sigma_{m,s}$. The multiplicity code $\mathrm{MRM}(q,m,d,s)$ is defined as follows. The alphabet of the code is $\Sigma_{m,s}$ and the length is $q^m$. Every polynomial $P(\boldsymbol X)\in\mathbb F_q[\boldsymbol X]$ of $\deg(P)\leq d$ defines a codeword by $(P^{(<s)}(\boldsymbol a))_{\boldsymbol a:\boldsymbol a\in\mathbb F_q^m}\in(\Sigma_{m,s})^{q^m}$.
We also let $\mathrm{MRS}(q,d,s):=\mathrm{Mult}(q,1,d,s)$ stand for Reed-Solomon multiplicity code.
The following lemma states the relationship between the derivatives of a polynomial to the derivatives of its restriction to a line.
\paragraph{Lemma 17 ([15], Sec 4).} Let $P\in\mathbb F[\boldsymbol X]$ be a multivariate polynomial where $\boldsymbol X=(X_1,\ldots,X_m)$. Let $\boldsymbol a,\boldsymbol b\in\mathbb F^m$ and define a univariate polynomial $P|_{L_{\boldsymbol a,\boldsymbol b}}(t)=P(\boldsymbol a t+\boldsymbol b)$. Then
\[(P|_{L_{\boldsymbol a,\boldsymbol b}})^{(j)}(t)=\sum_{\boldsymbol I:w(\boldsymbol I)=j}P^{(\boldsymbol I)}(\boldsymbol a t+\boldsymbol b)\cdot\boldsymbol a^{\boldsymbol I}.\]
\section*{3 Polynomials and tables (author context excerpts)}
A table is an element of $\Sigma_{m,s}^{q^m}$, i.e., a function mapping every evaluation point in $\mathbb F_q^m$ to an element in $\Sigma_{m,s}$. EVAL takes a multi-variate polynomial and returns its table of evaluations. More precisely,
\paragraph{Definition 39.} We define $\mathrm{EVAL}_{m,s}:\mathbb F_q[\boldsymbol X]\times\mathbb F_q^m\to\Sigma_{m,s}$ by
\[\mathrm{EVAL}_{m,s}(P;\boldsymbol a)=(P^{(\boldsymbol I)}(\boldsymbol a))_{w(\boldsymbol I)<s}.\]
Similarly, we define $\mathrm{EVAL}_{m,s}:\mathbb F_q[\boldsymbol X]\to(\Sigma_{m,s})^{q^m}$ by
\[\mathrm{EVAL}_{m,s}(P)=(\mathrm{EVAL}(P;\boldsymbol a))_{\boldsymbol a\in\mathbb F_q^m}.\]
We say $T\in(\Sigma_{m,s})^{q^m}$ is the table of $P\in\mathbb F_q[\boldsymbol X]$ if $\mathrm{EVAL}(P)=T$.
An element in $\Sigma_{m,s}$ contains information about all the derivatives of order up to $s$. Sometimes we would like to focus on a certain directional derivative. If $\sigma\in\Sigma_{m,s}$ and $\boldsymbol I\in\mathbb N^m$ with $w(\boldsymbol I)<s$, then $\sigma_{(\boldsymbol I)}\in\mathbb F_q$ is the entry of $\sigma$ that encodes the $\boldsymbol I$'th derivative. Similarly, if $f\in\Sigma_{m,s}^{q^m}$ is a table, i.e., $f:\mathbb F_q^m\to\Sigma_{m,s}$, then $f_{(\boldsymbol I)}:\mathbb F_q^m\to\mathbb F_q$ is defined by letting $f_{(\boldsymbol I)}(\boldsymbol x)$ be the $\boldsymbol I$'th entry of $f(\boldsymbol x)\in\Sigma_{m,s}$.
To understand the low-dimensional tests on tables, we need to define the restriction of tables to subspaces. If $T$ is a table $T:\mathbb F_q^m\to\Sigma_{m,s}$ and we want to restrict to the plane $\boldsymbol a t+\boldsymbol b r+\boldsymbol c$ the restriction $T_{\boldsymbol a,\boldsymbol b,\boldsymbol c}$ should be a table $\mathbb F_q^2\to\Sigma_{2,s}$. To this end we define the alphabet reduction map $\phi_{(\boldsymbol a,\boldsymbol b)}:\Sigma_{m,s}\to\Sigma_{2,s}$ by
\paragraph{Definition 48.}
\[(\phi_{(\boldsymbol a,\boldsymbol b)})(\boldsymbol z)_{\boldsymbol J=(j_1,j_2)}=\sum_{\boldsymbol I\in\mathbb N^m}z_{\boldsymbol I}\cdot\sum_{\substack{\boldsymbol I_1+\boldsymbol I_2=\boldsymbol I\\w(\boldsymbol I_1)=j_1,w(\boldsymbol I_2)=j_2}}\binom{\boldsymbol I}{\boldsymbol I_1}\boldsymbol a^{\boldsymbol I_1}\boldsymbol b^{\boldsymbol I_2}.\]
This map applies the ``chain rule'' to an element in $\Sigma_{m,s}$, in accordance with Lemma 18. We may then define
\[T_{\boldsymbol a,\boldsymbol b,\boldsymbol c}=\phi_{(\boldsymbol a,\boldsymbol b)}\circ T\circ\ell_{\boldsymbol a,\boldsymbol b,\boldsymbol c}.\]
Similarly, for $\boldsymbol h=(\boldsymbol h_0,\boldsymbol h_1,\ldots,\boldsymbol h_k)$ we define $\phi_{\boldsymbol h}$ and $T_{\boldsymbol h}$ by:
\[(\phi_{\boldsymbol h}(\boldsymbol z))_{\boldsymbol J}=\sum_{\boldsymbol I\in\mathbb N^m}z_{\boldsymbol I}\cdot\sum_{\substack{\boldsymbol I_1+\cdots+\boldsymbol I_k=\boldsymbol I\\w(\boldsymbol I_r)=j_r}}\binom{\boldsymbol I}{\boldsymbol I_1,\ldots,\boldsymbol I_k}\prod_{i=1}^k\boldsymbol h_k^{\boldsymbol I_k},\text{ and,}\]
\[T_{\boldsymbol h}(\boldsymbol Y)=\phi_{(\boldsymbol h)}\circ T\left(\boldsymbol h_0+\sum_{i=1}^kY_i\boldsymbol h_i\right).\]
\section*{5 Local testing above the RM testing dimension}
In this section we look at the local characterization and local testing of $\mathrm{MRM}(q,m,d,s)$ by dimension $k$ tests, when $k$ is above the RM testing dimension $t_{q,d}=\lceil(d+1)/(q-q/p)\rceil$ of $\mathrm{RM}(q,m,d)$. By Theorem 2 dimension $k$ subspaces give a local test (and hence also a local characterization) for $\mathrm{RM}(q,m,d)$. We show they also give a local test (and hence also a local characterization) for $\mathrm{MRM}(q,m,d,s)$ for $s<q$. There is, however, some parameter loss in the reduction as we next explain. To formally state the result we need some notation. For $\boldsymbol x\in\mathbb F_q^m$ let
\begin{align*}
H_k&=\{\boldsymbol h=(\boldsymbol h_0,\boldsymbol h_1,\ldots,\boldsymbol h_k)\mid\boldsymbol h_0,\ldots,\boldsymbol h_k\in\mathbb F_q^m,\dim(\operatorname{span}\{\boldsymbol h_1,\ldots,\boldsymbol h_k\})=k\},\text{ and,}\\
H_{k,\boldsymbol x}&=\{\boldsymbol h\in H_k\mid\boldsymbol x\in\boldsymbol h_0+\operatorname{Span}\{\boldsymbol h_1,\ldots,\boldsymbol h_k\}\}.
\end{align*}
By $\boldsymbol h\sim_kH$ (resp. $H_{k,\boldsymbol x}$) we mean a choosing $\boldsymbol h$ uniformly at random from $H_k$ (resp. $H_{k,\boldsymbol x}$). We let $A_{\boldsymbol h}=\boldsymbol h_0+\operatorname{Span}\{\boldsymbol h_1,\ldots,\boldsymbol h_k\}$ be the $k$-dimensional affine space defined by $\boldsymbol h$. We also recall the distance and rejection function from the introduction:
\paragraph{Definition 50.} Let $T:\mathbb F_q^m\to\Sigma_{m,s}$ be a table. Then
\begin{itemize}
\item $\delta(T,\mathrm{MRM}(q,m,d,s))$ is the distance between $T$ and the closest evaluation table of a degree $d$ polynomial, i.e., $\delta(T,\mathrm{MRM}(q,m,d,s))=\min_{G\in\mathrm{MRM}(q,m,d,s)}\{\delta(T,G)\}$, and,
\item $\mathrm{REJ}_{k,d}^{\mathrm{MRM}}(T)$ is the probability a dimension-$k$ test demonstrates that $T$ is not a degree $d$ polynomial.
\end{itemize}
With this notation we prove:
\paragraph{Theorem 51.} Let $\mathbb F_q$ be a field of size $q$, and assume $s\leq\min\{d,q-1\}$. Suppose for $\mathrm{RM}(q,m,d)$ there exists $\alpha>0$ and $c_0\leq1$ such that for every $f$
\[\mathrm{REJ}_{k,d}^{\mathrm{RM}}(f)\geq\min\{\alpha\cdot\delta(f,\mathrm{RM}(q,m,d)),c_0\}.\]
Then, for every $s<q$, for every $T$ we have
\[\mathrm{REJ}_{k,d}^{\mathrm{MRM}}(T)\geq\min\{\alpha'\cdot\delta(T,\mathrm{MRM}(q,m,d,s)),c_0\}\tag{6}\]
for
\[\alpha'=\left(1-\frac{s-1}{q}\right)\frac1{1+q^{d/(q-1)}\frac1\alpha}.\tag{7}\]
For example, assume $q$ is prime. If $d+1\leq q-1$ the testing dimension is $1$, meaning that lines are a good local test for $\mathrm{RM}(q,m,d)$. Then, Theorem 51 says that lines are also a good test for $\mathrm{MRM}(q,m,d,s)$, alas, with a larger coefficient. Similarly, if $q<d\leq2(q-1)$ the RM testing dimension is $2$ and therefore planes are a good local test for $\mathrm{RM}(q,m,d)$. Then, Theorem 51 says that planes are also a good test for $\mathrm{MRM}(q,m,d,s)$, with a larger coefficient. In the same vein 3-dimensional planes are a good local test for $\mathrm{MRM}(q,m,d,s)$ when $d\leq3(q-1)$.
To explain the intuition behind the theorem we first consider the characterization aspect of it. Suppose $k$ is above the RM testing dimension, and that the table $T$ passes all $\mathrm{MRM}(q,m,d,s)$ $k$-dimensional tests $\boldsymbol h$. Specifically, this means that for every $\boldsymbol h\in H_k$, the table $T$ restricted to the $k$-dimensional affine space $A_{\boldsymbol h}$ is consistent with a degree $d$ polynomial $P_{\boldsymbol h}$. Now, let $T_{(0)}$ be the table $T$ where at each entry we keep only the evaluation of the function itself and remove the evaluations that are associated with higher derivatives. Then, in particular, for every $\boldsymbol h\in H_k$, the table $T_{(0)}$ restricted to the $k$-dimensional space $A_{\boldsymbol h}$ is still consistent with the degree $d$ polynomial $P_{\boldsymbol h}$. As $k$ is above the RM testing dimension, there must be a unique degree $d$ polynomial $P$ that is consistent with $P_{\boldsymbol h}$ (and the table $T_{(0)}$) for every $\boldsymbol h\in H_k$. This $P$ is the only possible candidate for a low-degree explanation of the table $T$. What we need to check, and is indeed correct, is that since $T$ passes all dimension $k$ tests, $P$ is indeed consistent with $T$.
The testing case requires more technical details but is similar in spirit. We first, again, look at $T_{(0)}$ that contains only the function evaluations, and not the higher derivatives evaluations. If $T$ passes the test with high enough probability, then so does $T_{(0)}$, and this ensures the existence of a global degree-$d$ polynomial $P$ that agrees with most values of $T_{(0)}$. Again, what remains to be shown is that $P$ agrees with most values of $T$, which we indeed prove.
In short, one can informally say that Theorem 51 shows that the MRM testing dimension is not larger than the RM testing dimension, and that above the RM testing dimension one can get local testing for MRM codes. A natural question is whether the MRM testing dimension is equal to the RM testing dimension, or not. In Section 4 we saw that when the RM testing dimension is larger than $1$, so is the MRM testing dimension, and lines do not characterize the MRM code. One might be drawn to the conjecture that the RM and MRM testing dimensions coincide. However, surprisingly, in Section 6 we show that no matter what the RM testing dimension is, when $d<sq$ the MRM testing dimension is at most two (for a precise statement see Theorem 60). For example, for $\mathrm{MRM}(q,d,m,s)$ with $2q<d<3(q-1)$, the RM testing dimension is (roughly) $3$ while the MRM testing dimension is still $2$.
\subsection*{5.1 The proof}
\paragraph{Proof.} Let $T:\mathbb F_q^m\to\Sigma_{m,s}$ be a table. Let
\[\rho=\mathrm{REJ}_{k,d}^{\mathrm{MRM}}(T)\]
be the $k$-dimensional MRM test rejection probability. If $\rho\geq c_0$ we are done. Therefore, we assume $\rho<c_0$. We first utilize what we know at the zero level.
\paragraph{Claim 53.} Let $\delta^{\mathrm{RM}}=\delta_{q,m,d}^{\mathrm{RM}}$ be the distance of the $\mathrm{RM}(q,m,d)$ code. There exists a degree $d$ polynomial $P$ such that
\[\Pr_{\boldsymbol h\in H_k}((\phi_{\boldsymbol h}\circ T)|_{A_{\boldsymbol h}}\neq\mathrm{EVAL}(P|_{A_{\boldsymbol h}}))\stackrel{\mathrm{def}}=\varepsilon_0\leq\rho\left(1+\frac1{\alpha\cdot\delta^{\mathrm{RM}}}\right).\]
We call $\boldsymbol h$ good if $(\phi_{\boldsymbol h}\circ T)|_{A_{\boldsymbol h}}=\mathrm{EVAL}(P|_{A_{\boldsymbol h}})$ and bad otherwise. In this terminology, $\varepsilon_0=\Pr_{\boldsymbol h\in H_k}[\boldsymbol h\text{ is bad}]$.
\paragraph{Proof.} Let $T_{(0)}:\mathbb F_q^m\to\mathbb F_q$ be as in Section 3, i.e., the table where we keep only the entries of the function evaluations and ignore the evaluations of higher order derivatives. For a given affine space $A$ we have that $(T|_A)_{(0)}=T_{(0)}|_A$, and so if the former agrees with a degree $d$ polynomial, so does the latter. It follows that
\[\mathrm{REJ}_{k,d}^{\mathrm{RM}}(T_{(0)})\leq\rho.\]
By the hypothesis regarding $\mathrm{RM}(q,m,d)$ and using $\rho<c_0$ there exists a unique degree $d$ polynomial $P\in\mathbb F_q[\boldsymbol X]$ such that
\[\delta_{(0)}\stackrel{\mathrm{def}}=\delta(T_{(0)},P)\leq\rho/\alpha.\]
We notice that
\begin{align*}
\mathbb E_{\boldsymbol h\sim H_k}\delta(T_{(0)}|_{A_{\boldsymbol h}},P|_{A_{\boldsymbol h}})
&=\Pr_{\boldsymbol x\sim\mathbb F_q^m}(T_{(0)}(\boldsymbol x)\neq P(\boldsymbol x))\\
&=\delta(T_{(0)},P)=\delta_{(0)},
\end{align*}
because the subspaces in $A_{\boldsymbol h}$ for $\boldsymbol h\in H_k$, cover every point an equal number of times. Therefore, by Markov's inequality,
\[\Pr_{\boldsymbol h\sim H_k}(\delta(T_{(0)}|_{A_{\boldsymbol h}},P|_{A_{\boldsymbol h}})\geq\delta^{\mathrm{RM}})\leq\frac{\delta_{(0)}}{\delta^{\mathrm{RM}}}.\tag{8}\]
Also, by assumption,
\[\Pr_{\boldsymbol h\in H_k}((\phi_{\boldsymbol h}\circ T)|_{A_{\boldsymbol h}}\notin\mathrm{MRM}(q,k,d,s))\leq\rho.\tag{9}\]
This means that except for probability $\rho$ over $\boldsymbol h$, $(\phi_{\boldsymbol h}\circ T)|_{A_{\boldsymbol h}}$ agrees with $\mathrm{EVAL}(P_{\boldsymbol h})$ for some $k$-variate degree $d$ polynomial $P_{\boldsymbol h}$. When this happens it also holds that $T_{(0)}$ agrees with $P_{\boldsymbol h}$ over $A_{\boldsymbol h}$ (because $\phi_{\boldsymbol h}$ does not change the zero level).
Thus, Equations (8) and (9) together imply that except for probability $\rho+\delta_0/\delta^{\mathrm{RM}}\leq\rho(1+1/(\alpha\delta^{\mathrm{RM}}))$ over $\boldsymbol h$, we simultaneously have that $T_{(0)}=P_{\boldsymbol h}$ over $A_{\boldsymbol h}$ and that $\delta(T_{(0)}|_{A_{\boldsymbol h}},P|_{A_{\boldsymbol h}})<\delta^{\mathrm{RM}}$. When both events happen we conclude that
\[\delta(P_{\boldsymbol h},P|_{A_{\boldsymbol h}})<\delta^{\mathrm{RM}}.\]
As $P_{\boldsymbol h}$ and $P|_{A_{\boldsymbol h}}$ are degree $d$ polynomials on $A_{\boldsymbol h}$ and are closer than $\delta^{\mathrm{RM}}$, it must be the that in fact $P_{\boldsymbol h}=P|_{A_{\boldsymbol h}}$. Thus, $(\phi_{\boldsymbol h}\circ T)|_{A_{\boldsymbol h}}$ is a valid $\mathrm{MRM}(q,k,d,s)$ table, and it is the table of the polynomial $P|_{A_{\boldsymbol h}}$, i.e., $(\phi_{\boldsymbol h}\circ T)|_{A_{\boldsymbol h}}=\mathrm{EVAL}(P|_{A_{\boldsymbol h}})$ as desired.\hfill$\triangleleft$
Our goal is to bound $\delta(\mathrm{EVAL}(P),T)$. This means that at most $\boldsymbol x\in\mathbb F_q^m$ we should have $T(\boldsymbol x)=\mathrm{EVAL}(P;\boldsymbol x)$. $T(\boldsymbol x)$ and $\mathrm{EVAL}(P;\boldsymbol x)$ contains values for all the $m$-variate directional derivatives of order up to $s$. Our handle on these values is Claim 53, that shows that most $\boldsymbol h$ are good, meaning that the $k$-variate derivatives of $P|_{A_{\boldsymbol h}}$ and $\phi_{\boldsymbol h}\circ T|_{A_{\boldsymbol h}}$ are the same. We notice that every such $k$-variate derivative is a linear combination (dependent on $\boldsymbol h$) of the $m$-variate derivatives. If $\boldsymbol x$ is such that for many $\boldsymbol h\in H_{k,\boldsymbol x}$, $\boldsymbol h$ is good, then for that $\boldsymbol x$ we get many linear equations on the $m$-variate derivatives. Our task is to prove that for many $\boldsymbol x$ there are enough good $\boldsymbol h\in H_{k,\boldsymbol x}$ to force the underlying $m$-variate derivatives of $P$ and $T$ to agree.
\paragraph{Claim 54.} We say $\boldsymbol x\in\mathbb F_q^m$ is bad if $\Pr_{\boldsymbol h\in H_{k,\boldsymbol x}}[\boldsymbol h\text{ is bad}]\geq1-(s-1)/q$ and good otherwise. Then
\[\Pr_{\boldsymbol x\in\mathbb F_q^m}[\boldsymbol x\text{ is bad}]\leq\frac q{q-(s-1)}\cdot\Pr_{\boldsymbol h\in H_k}[\boldsymbol h\text{ is bad}].\]
\paragraph{Proof.} We have
\[\mathbb E_{\boldsymbol x\in\mathbb F_q^m}\left[\Pr_{\boldsymbol h\in H_{k,\boldsymbol x}}(\boldsymbol h\text{ is bad})\right]=\Pr_{\boldsymbol h\in H_k}(\boldsymbol h\text{ is bad})=\varepsilon_0,\]
where $\varepsilon_0$ is as in Claim 53. This is because choosing a uniform $\boldsymbol h\in H_k$ is the same as first choosing a uniform $\boldsymbol x\in\mathbb F_q^m$ and then choosing a uniform $\boldsymbol h\in H_{k,\boldsymbol x}$. Therefore, for every $c>1$, by Markov's inequality,
\[\Pr_{\boldsymbol x\in\mathbb F_q^m}\left[\Pr_{\boldsymbol h\in H_{k,\boldsymbol x}}(\boldsymbol h\text{ is bad})\geq c\cdot\varepsilon_0\right]\leq1/c.\]
Choosing $c=(q-(s-1))/(q\cdot\varepsilon_0)$ gives the result.\hfill$\triangleleft$
We note that if most $k$-dimensional subspaces are good, then most lines are:
\paragraph{Claim 55.} For every $\boldsymbol x\in\mathbb F_q^m$ and $k\geq1$,
\[\Pr_{\boldsymbol u\in H_{1,\boldsymbol x}}[\boldsymbol u\text{ is bad}]\leq\Pr_{\boldsymbol h\in H_{k,\boldsymbol x}}[\boldsymbol h\text{ is bad}].\]
\paragraph{Proof.} Fix $\boldsymbol x$. For every $\boldsymbol h\in H_{k,\boldsymbol x}$, If $\boldsymbol h$ is good, then $(\phi_{\boldsymbol h}\circ T)|_{A_{\boldsymbol h}}=\mathrm{EVAL}(P|_{A_{\boldsymbol h}})$. This implies, in particular, that for every $\boldsymbol u\in H_{1,\boldsymbol x}$ such that $A_{\boldsymbol u}\subseteq A_{\boldsymbol h}$, we also have that $(\phi_{\boldsymbol u}\circ T)|_{A_{\boldsymbol u}}=\mathrm{EVAL}(P|_{A_{\boldsymbol u}})$. The result then follows because we can sample $\boldsymbol u\in H_{1,\boldsymbol x}$ by first sampling $\boldsymbol h\in A_{k,\boldsymbol x}$ and then choosing a random $\boldsymbol u\in H_{1,\boldsymbol x}$ such that $A_{\boldsymbol u}\subseteq A_{\boldsymbol h}$.\hfill$\triangleleft$
We now fix any good $\boldsymbol x$. We will prove:
\paragraph{Claim 56.} Suppose $\boldsymbol x\in\mathbb F_q^m$ is good. Then for any $m$-variate direction $\boldsymbol I$ with $w(\boldsymbol I)<s$ we have
\[T_{(\boldsymbol I)}(\boldsymbol x)=P^{(\boldsymbol I)}(\boldsymbol x).\tag{10}\]
Once we prove the claim we can conclude the proof of the theorem because:
\begin{align*}
\delta(T,\mathrm{EVAL}(P))&\leq\Pr_{\boldsymbol x\in\mathbb F_q^m}(\boldsymbol x\text{ is bad})\leq\frac q{q-(s-1)}\cdot\varepsilon_0\\
&\leq\frac q{q-(s-1)}\cdot\rho\cdot\left(1+\frac1{\alpha\cdot\delta^{\mathrm{RM}}}\right)\\
&\leq\rho\left(1-\frac q{s-1}\right)^{-1}\left(1+\frac{q^{d/(q-1)}}\alpha\right).
\end{align*}
\paragraph{Proof of Claim 56.} Fix a good $\boldsymbol x\in\mathbb F_q^m$ and let $w_0<s$. We will show Equation (10) simultaneously for all $m$-variate directions $\boldsymbol I$ with $w(\boldsymbol I)=w_0$. We know that
\[\Pr_{\boldsymbol u\in H_{1,\boldsymbol x}}(\boldsymbol u\text{ is bad})\leq\Pr_{\boldsymbol h\in H_{k,\boldsymbol x}}(\boldsymbol h\text{ is bad})<1-\frac{s-1}q,\tag{11}\]
where the first inequality is by Claim 55 and the second because $\boldsymbol x$ is good.
Suppose $\boldsymbol u=(\boldsymbol b,\boldsymbol a)$ is good (where $\boldsymbol b,\boldsymbol a\in\mathbb F_q^m$). Then, by definition, $(\phi_{\boldsymbol u}\circ T)|_{A_{\boldsymbol u}}=\mathrm{EVAL}(P|_{A_{\boldsymbol u}})$. In particular $\phi_{\boldsymbol u}(T(\boldsymbol x))=\mathrm{EVAL}(P|_{A_{\boldsymbol u}};\boldsymbol x)$, because $\boldsymbol x\in A_\ell$. Now:
\begin{itemize}
\item By Definition 48 (and plugging $k=1$)
\[(\phi_{\boldsymbol u}(T(\boldsymbol x)))_{w_0}=\sum_{\boldsymbol I,w(\boldsymbol I)=w_0}T(\boldsymbol x)_{(\boldsymbol I)}\cdot\boldsymbol a^{\boldsymbol I},\]
\item Also, by Lemma 17 (plugging $t=0$),
\[(\mathrm{EVAL}(P|_{A_{\boldsymbol u}};\boldsymbol x))_{w_0}=\sum_{\boldsymbol I,w(\boldsymbol I)=w_0}P^{\boldsymbol I}(\boldsymbol x)\cdot\boldsymbol a^{\boldsymbol I}.\]
\end{itemize}
Thus, a good $\boldsymbol u=(\boldsymbol b,\boldsymbol a)$ gives the liner equation
\[\sum_{\boldsymbol I,w(\boldsymbol I)=w_0}(P^{\boldsymbol I}(\boldsymbol x)-T(\boldsymbol x)_{(\boldsymbol I)})\cdot\boldsymbol a^{\boldsymbol I}=0,\]
where the variables are $v_{\boldsymbol I}=P^{\boldsymbol I}(\boldsymbol x)-T(\boldsymbol x)_{(\boldsymbol I)}$ for every $m$-variate direction $\boldsymbol I$ of total weight exactly $w_0$. Thus, Equation (11) implies that
\[\Pr_{\boldsymbol a\in\mathbb F_q^m\setminus\{\boldsymbol0\}}\left(\sum_{\boldsymbol I,w(\boldsymbol I)=w_0}(T(\boldsymbol x)_{(\boldsymbol I)}-P^{(\boldsymbol I)}(\boldsymbol x))\cdot\boldsymbol a^{\boldsymbol I}=0\right)>\frac{s-1}q.\]
Now, look at the polynomial $f_{\boldsymbol x}\in\mathbb F_q[X_1,\ldots,X_m]$ defined by
\[f_{\boldsymbol x}(\boldsymbol a)=\sum_{\boldsymbol I,w(\boldsymbol I)=w_0}(T(\boldsymbol x)_{(\boldsymbol I)}-P^{(\boldsymbol I)}(\boldsymbol x))\cdot\boldsymbol a^{\boldsymbol I}.\]
$f_{\boldsymbol x}$ is an $m$-variate, degree $w_0$ homogeneous polynomial, and it is $0$ with probability larger than $(s-1)/q\geq w_0/q=\deg(f_{\boldsymbol x})/q$. By the Schwartz-Zippel lemma, it must be the $0$ polynomial. Therefore, $T(\boldsymbol x)_{(\boldsymbol I)}-P^{(\boldsymbol I)}(\boldsymbol x)$ for all $\boldsymbol I$ with $w(\boldsymbol I)=w_0$ as desired.\hfill$\triangleleft$\par\hfill$\triangleleft$
\paragraph{Remark 57.} Another way to view the argument, is that each $\boldsymbol a\in\mathbb F_q^m$ is an evaluation point of a homogeneous $\mathrm{RM}(q,m,w_0)$ codeword, and therefore each good $\boldsymbol u=(\boldsymbol b,\boldsymbol a)$ gives a zero coordinate of the codeword. If the number of good $\boldsymbol u$ is too large, we get too many zeroes, and therefore the codeword must be the zero codeword, meaning that the values of the variables are zero as we wish.
\paragraph{Remark 58.} The argument we use has the information that many $k$-dimensional restrictions are good, but then chooses to reduce this knowledge to the weaker statement that for many $\boldsymbol x$, for many lines passing through $\boldsymbol x$, the linear restrictions are good. It seems that using the stronger statement might give a better code and improve the parameters, but we have not succeeded yet in analyzing this.

% Additional faithful author context for the main plane tester, official Sections2–3 and Corollary52.
\section*{Additional author algebraic context}
Recall The definition of the binomial and multinomial coefficients for natural numbers:
\[\binom ab\stackrel{\mathrm{def}}=\frac{a!}{b!(a-b)!},\qquad\binom a{b_1,\ldots,b_\ell}=\frac{a!}{b_1!\cdots b_\ell!},\]
where $\sum b_i=a$. We extend this definition to $\boldsymbol I,\boldsymbol I_1,\boldsymbol J,\boldsymbol J_1,\ldots,\boldsymbol J_\ell\in\mathbb N^m$ where $\boldsymbol J=\boldsymbol J_1+\ldots+\boldsymbol J_\ell$ by
\[\binom{\boldsymbol I}{\boldsymbol I_1}\stackrel{\mathrm{def}}=\prod_{k=1}^m\binom{I_k}{(I_1)_k},\qquad\binom{\boldsymbol J}{\boldsymbol J_1,\ldots,\boldsymbol J_\ell}\stackrel{\mathrm{def}}=\prod_{k=1}^m\binom{J_k}{(J_1)_k,\cdots,(J_\ell)_k}.\]
We mention Lucas theorem, that if $p$ is prime, $q=p^w$ for $w\in\mathbb N$, and $a,b\in\mathbb N$ have base $p$ representation $a=\sum_{\ell=0}^wa_\ell p^\ell$, $b=\sum_{\ell=0}^wb_\ell p^\ell$ with $0\leq a_\ell,b_\ell<p$, then
\[\binom ab\bmod p=\prod_{\ell=0}^w\binom{a_\ell}{b_\ell},\tag{4}\]
where we use the convention that $\binom cd=0$ when $d>c$. Thus, $\binom ab\bmod p\neq0$ iff $a_\ell\geq b_\ell$ for all $\ell=0,\ldots,w-1$.
Finally, we let $g\in\mathbb F_q[X]$ denote the polynomial $g(X)=X^q-X$. For $\boldsymbol X=(X_1,\ldots,X_m)$ and $\boldsymbol I=(i_1,\ldots,i_m)\in\mathbb N^m$ we let $g(\boldsymbol X)^{\boldsymbol I}$ denote $\prod_{k=1}^m(g(X_k))^{i_k}$.
\paragraph{Lemma 18.} Let $P\in\mathbb F[\boldsymbol X]$ be a multivariate polynomial where $\boldsymbol X=(X_1,\ldots,X_m)$. Let $\boldsymbol a,\boldsymbol b\in\mathbb F^m$ and define a bivariate polynomial by $P|_{P_{\boldsymbol a,\boldsymbol b,\boldsymbol c}}(t,r)=P(\boldsymbol at+\boldsymbol br+\boldsymbol c)$. Then for $(j_1,j_2)\in\mathbb N^2$:
\[(P|_{P_{\boldsymbol a,\boldsymbol b,\boldsymbol c}})^{(j_1,j_2)}(t,r)=\sum_{\boldsymbol I\in\mathbb N^m}P^{(\boldsymbol I)}(\boldsymbol at+\boldsymbol br+\boldsymbol c)\cdot\sum_{\substack{\boldsymbol I_1+\boldsymbol I_2=\boldsymbol I\\w(\boldsymbol I_1)=j_1,w(\boldsymbol I_2)=j_2}}\binom{\boldsymbol I}{\boldsymbol I_1}\boldsymbol a^{\boldsymbol I_1}\boldsymbol b^{\boldsymbol I_2}.\]
\paragraph{Proof.} Given $R_1,R_2\in\mathbb F$, we write the expression $P(\boldsymbol a(t+R_1)+\boldsymbol b(r+R_2)+\boldsymbol c)$ in two different ways. On the one hand,
\begin{align*}
P(\boldsymbol a(t+R_1)+\boldsymbol b(r+R_2)+\boldsymbol c)
&=P|_{P_{\boldsymbol a,\boldsymbol b,\boldsymbol c}}(t+R_1,r+R_2)\\
&=P|_{P_{\boldsymbol a,\boldsymbol b,\boldsymbol c}}((t,r)+(R_1,R_2))\\
&=\sum_{j_1,j_2\in\mathbb N}(P|_{P_{\boldsymbol a,\boldsymbol b,\boldsymbol c}})^{(j_1,j_2)}(t,r)R_1^{j_1}R_2^{j_2}.
\end{align*}
On the other hand,
\begin{align*}
P(\boldsymbol a(t+R_1)+\boldsymbol b(r+R_2)+\boldsymbol c)
&=P(\boldsymbol at+\boldsymbol br+\boldsymbol c+R_1\boldsymbol a+R_2\boldsymbol b)\\
&=\sum_{\boldsymbol I\in\mathbb N^m}P^{(\boldsymbol I)}(\boldsymbol at+\boldsymbol br+\boldsymbol c)\cdot(R_1\boldsymbol a+R_2\boldsymbol b)^{\boldsymbol I}\\
&=\sum_{\boldsymbol I\in\mathbb N^m}P^{(\boldsymbol I)}(\boldsymbol at+\boldsymbol br+\boldsymbol c)\cdot\prod_{k=1}^m(a_kR_1+b_kR_2)^{i_k}\\
&=\sum_{\boldsymbol I\in\mathbb N^m}P^{(\boldsymbol I)}(\boldsymbol at+\boldsymbol br+\boldsymbol c)
\sum_{\substack{\boldsymbol I_1+\boldsymbol I_2=\boldsymbol I\\w(\boldsymbol I_1)=j_1\\w(\boldsymbol I_2)=j_2}}\binom{\boldsymbol I}{\boldsymbol I_1}\boldsymbol a^{\boldsymbol I_1}\boldsymbol b^{\boldsymbol I_2}R_1^{j_1}R_2^{j_2}.
\end{align*}
Comparing coefficients of $R_1^{j_1}R_2^{j_2}$ for every $\boldsymbol J=(j_1,j_2)\in\mathbb N^2$ we get the result.\hfill$\triangleleft$
We will also be interested in the restrictions of polynomials to general $k$-dimensional subspaces. Let $P|_{P_{\boldsymbol h_0,\boldsymbol h_1,\ldots,\boldsymbol h_k}}(\boldsymbol Y)=P(\boldsymbol h_0+\sum_{i=1}^k\boldsymbol h_iY_i)$. Then, similarly to Lemma 18,
\paragraph{Lemma 19.}
\[(P|_{P_{\boldsymbol h_0,\boldsymbol h_1,\ldots,\boldsymbol h_k}})^{\boldsymbol J}(\boldsymbol Y)=\sum_{\boldsymbol I\in\mathbb N^m}P^{(\boldsymbol I)}\left(\boldsymbol h_0+\sum_{i=1}^k\boldsymbol h_iY_i\right)\cdot\sum_{\substack{\boldsymbol I_1+\cdots+\boldsymbol I_k=\boldsymbol I\\w(\boldsymbol I_r)=j_r}}\binom{\boldsymbol I}{\boldsymbol I_1,\ldots,\boldsymbol I_k}\prod_{i=1}^k\boldsymbol h_k^{\boldsymbol I_k}.\]
The proof is identical to the proof of Lemma 18.
\subsection*{2.3 Grobner bases and Nullstellensatz}
The theory of Grobner bases describes the structure of ideals in the ring $R=\mathbb F[\boldsymbol X]$ and we briefly explain some of the essential concepts of this theory. We refer to [5] for a thorough treatment of this theory.
\paragraph{Definition 20.} A monomial order $\succ$ on $R$ is a relation $\succ$ on $\mathbb Z_{\geq0}^n$, or equivalently a relation on the set of monomials $x^\alpha$, $\alpha\in\mathbb Z_{\geq0}^n$ satisfying:
\begin{enumerate}
\item $\succ$ is a total ordering.
\item If $\alpha\succ\beta$ and $\gamma\in\mathbb Z_{\geq0}^n$ then $\alpha+\gamma\succ\beta+\gamma$.
\item $\succ$ is a well-ordering. I.e, every non-empty $A\subset\mathbb Z_{\geq0}^n$ has a minimal element.
\end{enumerate}
\paragraph{Example 21 (Lexicographic order).} Let $\alpha,\beta\in\mathbb Z_{\geq0}^m$. We say $\alpha\succ_{\mathrm{lex}}\beta$ if the minimal $i$ which satisfies $\alpha_i\neq\beta_i$, also satisfies $\alpha_i>\beta_i$.
\paragraph{Example 22 (Total degree lexicographic order).} The total degree lexicographic order is defined as follows: A monomial $m_1$ is greater than $m_2$ if it has higher total degree, where ties are broken lexicographically (i.e $X_1>X_2>\cdots>X_m$). More formally, let $\alpha,\beta\in\mathbb Z_{\geq0}^m$. Then $\alpha\succ_{\mathrm{tot}}\beta$ if $w(\alpha)=\sum\alpha_i>w(\beta)=\sum\beta_i$ or, $w(\alpha)=w(\beta)$ and $\alpha\succ_{\mathrm{lex}}\beta$.
\paragraph{Definition 23.} Let $f(\boldsymbol X)=\sum_{\boldsymbol I}a_{\boldsymbol I}\boldsymbol X^{\boldsymbol I}$ and $\succ$ a monomial order.
\begin{enumerate}
\item The multi-degree of $f$ is $\operatorname{multideg}(f)=\max\{\boldsymbol I\mid a_{\boldsymbol I}\neq0\}$ where the maximum is taken w.r.t $\succ$.
\item The leading coefficient of $f$ is $\mathrm{LC}(f)=a_{\operatorname{multideg}(f)}\in\mathbb F$.
\item The leading monomial of $f$ is $\mathrm{LM}(f)=\boldsymbol X^{\operatorname{multideg}(f)}$.
\item The leading term of $f$ is $\mathrm{LT}(f)=\mathrm{LC}(f)\cdot\mathrm{LM}(f)$.
\end{enumerate}
\paragraph{Definition 24 (Multivariate polynomial division).} Let $\succ$ be a monomial order on $\mathbb Z_{\geq0}^m$, and let $F=\{f_1,\ldots,f_k\}$ be a set of $k$ polynomials in $\mathbb F[\boldsymbol X]$. Then every $f\in\mathbb F[\boldsymbol X]$ can be written as
\[f=q_1f_1+\cdots+q_kf_k+r,\]
where $q_i,r\in\mathbb F[\boldsymbol X]$, and either $r=0$ or $r$ is a linear combination, with coefficients in $\mathbb F$, of monomials, none of which is divisible by any of $\mathrm{LT}(f_1),\ldots,\mathrm{LT}(f_k)$. We call $r$ a remainder of the division by $F$. Moreover, $\operatorname{multideg}(q_if_i)\leq\operatorname{multideg}(f)$ for every $i\in[s]$. The remainder $r$ is not necessarily unique, and might depend on the order of division.
\paragraph{Definition 25.} Let $\{0\}\neq I\subseteq\mathbb F[\boldsymbol X]$ be an ideal. Fix a monomial ordering on $\mathbb F[\boldsymbol X]$. Then
\begin{enumerate}
\item We denote by $\mathrm{LT}(I)$ the set of leading terms of non-zero elements in $I$.
\[\mathrm{LT}(I)=\{\mathrm{LT}(f)\mid f\in I\setminus\{0\}\}.\]
\item We denote by $\langle\mathrm{LT}(I)\rangle$ the ideal generated by the elements in $\mathrm{LT}(I)$.
\end{enumerate}
\paragraph{Definition 26.} Let $\{0\}\neq I\subseteq\mathbb F[\boldsymbol X]$ be an ideal. Fix a monomial order on $\mathbb F[\boldsymbol X]$. A subset $G=\{g_1,\ldots,g_t\}\subset I$ is said to be a Grobner basis for $I$, if
\[\langle\mathrm{LT}(g_1),\ldots,\mathrm{LT}(g_t)\rangle=\langle\mathrm{LT}(I)\rangle.\]
\paragraph{Note 27.} Every ideal $I\subset\mathbb F[X_1,\ldots,X_m]$ is finitely generated and has a Grobner basis.
The importance of a Grobner basis, is that it gives a natural way of choosing representatives for the quotient space $\mathbb F[X_1,\ldots,X_m]/I$.
\paragraph{Theorem 28 ([5], Section 2, Proposition 1).} Let $I\subset\mathbb F[X_1,\ldots,X_m]$ be an ideal and $G=\{g_1,\ldots,g_k\}$ a Grobner basis. Then given $f\in\mathbb F[\boldsymbol X]$ there is a unique $r\in\mathbb F[\boldsymbol X]$ such that there is a $g\in I$ such that $f=g+r$ and no term of $r$ is divisible by any of $\mathrm{LT}(g_1),\ldots,\mathrm{LT}(g_k)$. We call this $r$, the reduced form of $f$ (relative to $I$).
Note that the reduced form of any polynomial is equivalent to this polynomial modulo $I$. Thus, the theorem gives us a natural way of choosing representatives modulo $I$.
\paragraph{Theorem 29.} Let $R=\mathbb F[\boldsymbol X]$ be the ring of polynomials, and $I\subset R$ an ideal. Let $G$ be a Grobner basis for $I$. Then the set
\[B=\{M(\boldsymbol X)\mid M\text{ is a monomial not divisible by any }\mathrm{LT}(g)\text{ for }g\in G\},\]
is a basis for $R/I$.
The following criterion determines whether $G$ is a Grobner basis.
\paragraph{Definition 30 (LCM and S polynomials).} Let $f,g\in\mathbb F[\boldsymbol X]$ be non-zero polynomials. Let $\alpha=\operatorname{multideg}(f)$ and $\beta=\operatorname{multideg}(g)$.
\begin{enumerate}
\item The least common multiple of $\mathrm{LM}(f)$ and $\mathrm{LM}(g)$, denoted $\mathrm{LCM}(\mathrm{LM}(f),\mathrm{LM}(g))$, is $\boldsymbol X^\gamma$, where $\gamma=(\gamma_1,\ldots,\gamma_m)$ and $\gamma_i=\max\{\alpha_i,\beta_i\}$ for each $i$.
\item The $S$-polynomial of $f$ and $g$ is
\[S(f,g)=\frac{\mathrm{LCM}(\mathrm{LM}(f),\mathrm{LM}(g))}{\mathrm{LT}(f)}\cdot f-\frac{\mathrm{LCM}(\mathrm{LM}(f),\mathrm{LM}(g))}{\mathrm{LT}(g)}\cdot g.\]
\end{enumerate}
\paragraph{Theorem 31 (Buchberger's Criterion, [5], Sec 6).} Let $I\subset\mathbb F[\boldsymbol X]$ be an ideal. Then a basis $G=\{g_1,\ldots,g_k\}$ of $I$ is a Grobner basis of $I$ if and only if for all pairs $i\neq j$, the remainder on division of $S(g_i,g_j)$ by $G$ is zero.
Note that we always have $S=S(g_i,g_j)\in I$ by the definition of $S$. When saying the remainder of the division by $G$ is zero, we mean that there are $\{f_i\}$, such that $S=\sum f_ig_i$ and $\operatorname{multideg}(f_ig_i)\leq\operatorname{multideg}(S)$ for every $i$ (as in Definition 24).
\paragraph{Theorem 32 (Combinatorial Nullstellensatz, [1]).} Let $\mathbb F$ be a field, and $A_1,\ldots,A_m\subseteq\mathbb F$. Let $g_i(X)=\prod_{\alpha\in A_i}(X-\alpha)$ for $i=1,\ldots,m$. Assume a polynomial $f\in\mathbb F[\boldsymbol X]$ satisfies $f(\alpha)=0$ for all $\alpha\in A_1\times\cdots\times A_m$. Then there are $h_1,\ldots,h_t$ such that
\[f=\sum h_ig_i,\]
and $\deg(h_i)+\deg(g_i)\leq\deg(f)$ for all $i$.
When $A_i=\mathbb F_q$ denote
\[g(X)=\prod_{\alpha\in\mathbb F_q}(X-\alpha)=X^q-X.\]
Also, let $I_m$ denote the ideal $I_m=\{f\in\mathbb F_q[\boldsymbol X]\mid\forall\alpha\in\mathbb F_q^m\ f(\alpha)=0\}$.
\paragraph{Corollary 33.} $I_m=\langle g(X_1),\ldots,g(X_m)\rangle$.
\paragraph{Proof.} Let $f\in I_m$. By Theorem 32, taking $S_i=\mathbb F_q$ for every $i$, we get that $f=\sum h_ig_i$ for some $\{h_i\}$ and so $f\in\langle g_k\rangle_{k\in[m]}$. The other inclusion is trivial since $g(X_i)=X_i^q-X_i$ vanishes on $\mathbb F_q^m$ for every $i$.\hfill$\triangleleft$
Using Theorem 31 it can be easily proved that:
\paragraph{Claim 34.} $G=\{g(X_k)\}_{k=1}^m$ is a Grobner basis for $I_m$ relative to the total degree lexicographic order.
Let $s\in\mathbb N$ and let $I_{m,s}$ denote the ideal
\[I_{m,s}=\{f\in\mathbb F_q[\boldsymbol X]\mid\forall\alpha\in\mathbb F_q^m\ \mathrm{Mult}(f;\alpha)\geq s\}.\]
In this notation, $I_{m,1}=I_m$ defined before. For every $\boldsymbol I=(I_1,\ldots,I_m)\in\mathbb N^m$ define
\[g(\boldsymbol X)^{\boldsymbol I}=\prod_{k=1}^mg(X_k)^{I_k}.\]
\paragraph{Theorem 35 (Combinatorial Nullstellensatz with multiplicity, [3], Sec 3).} $I_{m,s}=\langle g(\boldsymbol X)^{\boldsymbol I}\rangle_{w(\boldsymbol I)=s}$. Furthermore, the set $G_{m,s}=\{g(\boldsymbol X)^{\boldsymbol I}\}_{w(\boldsymbol I)=s}$ is a Grobner basis for $I_{m,s}$.
\paragraph{Proof.} To see that $I_{m,s}$ is indeed an ideal, fix $f\in I_{m,s}$ and $h\in\mathbb F[\boldsymbol X]$. Then for $r<s$: $(hf)^{(r)}=\sum_{i=0}^rf^{(i)}\cdot h^{(r-i)}=0$ and so $hf\in I_{m,s}$. Also, clearly, $g(\boldsymbol X)^{\boldsymbol I}\in I_{m,s}$ for every $\boldsymbol I$ with $w(\boldsymbol I)=s$. We need to show $G_{m,s}$ is a Grobner basis for $I_{m,s}$,
[3, Section 3] show $G_{m,s}$ generates $I_{m,s}$ and, furthermore, $f=\sum_{b:w(\boldsymbol I)=s}g(\boldsymbol X)^{\boldsymbol I}h_{\boldsymbol I}$ for some $h_{\boldsymbol I}$ with $\deg(h_{\boldsymbol I})\leq\deg(f)-s\deg(g)$. In particular, this is true for the $S$ polynomials in Theorem 31. I.e, every such $S$ polynomial can be expressed as $S=\sum g(\boldsymbol X)^{\boldsymbol I}h_{\boldsymbol I}$ where $\deg(g(\boldsymbol X)^{\boldsymbol I}h_{\boldsymbol I})\leq\deg(S)$. By Buchberger's criterion $g(\boldsymbol X)^{\boldsymbol I}$ is a Grobner basis.\hfill$\triangleleft$
Finally, we look at equality modulo $I_{m,1}$.
\paragraph{Definition 36.} For $n_1,n_2\in\mathbb N$ we say $n_1=_{\mathbb F_q}n_2$ iff $x^{n_1}=x^{n_2}\bmod I_{1,1}$. Equivalently, $n_1=_{\mathbb F_q}n_2$ iff
\begin{itemize}
\item $n_1=n_2$, or,
\item $\min(n_1,n_2)>0$ and $n_1=n_2\bmod(q-1)$.
\end{itemize}
\paragraph{Definition 37.} Let $\boldsymbol A,\boldsymbol B\in\mathbb N^m$. We say $\boldsymbol A=_{\mathbb F_q}\boldsymbol B$ for iff $A_k=_{\mathbb F_q}B_k$ for every $1\leq k\leq m$.
We also record:
\paragraph{Claim 38.} Let $a,b\in\mathbb N$. If $a=_{\mathbb F_q}b$ and $a<q$ then $a\leq b$.
\paragraph{Proof.} If $b<q$ then $a=_{\mathbb F_q}b$ implies $a=b$. Otherwise $a<q\leq b$.\hfill$\triangleleft$
\section*{3 Polynomials and tables (continued author context)}
Note that every polynomial determines its table $\mathrm{EVAL}(P)$. However, two different polynomials might have the same table if their difference is the zero table on $\mathbb F_q^m$. From Theorem 35,
\[\operatorname{Ker}(\mathrm{EVAL}_{m,s})=I_{m,s}=\langle g(\boldsymbol X)^{\boldsymbol I}\rangle_{w(\boldsymbol I)=s}.\]
Since $I_{m,s}$ does not contain any non-zero polynomial of total degree less than $sq$, we conclude that:
\paragraph{Corollary 40.} $\mathrm{EVAL}_{m,s}$ is injective on polynomials of total degree less than $sq$.
\paragraph{Claim 41.} $\mathrm{EVAL}_{m,s}$ is onto $\Sigma_{m,s}^{q^m}$.
\paragraph{Proof.} We will show that
\[\dim(\mathbb F_q[X_1,\ldots,X_m]/I_{m,s})\geq\dim((\Sigma_{m,s})^{q^m})=q^m\cdot\binom{m+s-1}{m},\]
where the dimension is over $\mathbb F_q$. This implies that the image of $\mathrm{EVAL}_{m,s}$ is everything, and the mapping is injective.
To see that consider the set
\[B_{m,s}=\left\{g(\boldsymbol X)^{\boldsymbol I}\cdot\prod_{k=1}^mX_k^{j_k}\mid w(\boldsymbol I)<s,0\leq j_k<q\right\}.\tag{5}\]
The elements in $B_{m,s}$ have different multi degree and therefore are independent. Also elements in $B_{m,s}$ are monomials of degree smaller than $sq$, and therefore are $I_{m,s}$ reduced. Thus,
\begin{align*}
\dim(\mathbb F_q[\boldsymbol X]/I_{m,s})&\geq\dim\operatorname{Span}(B_{m,s})\\
&=|\{(\boldsymbol I,\boldsymbol J)\in\mathbb N^m\times\mathbb N^m\mid w(\boldsymbol I)<s,0\leq j_k<q\}|\\
&=\binom{m+s-1}m\cdot q^m,
\end{align*}
as desired.\hfill$\triangleleft$
\paragraph{Corollary 42.} The set $B_{m,s}$ is a basis for $\mathbb F[\boldsymbol X]/I_{m,s}$.
\paragraph{Example 43.} When $m=1$ the set $\{g(x)^ix^j\}_{i\in\mathbb N,j<q}$ is a basis of $\mathbb F[\boldsymbol X]$. Thus, for every $\ell\in\mathbb N$ there are $\beta_{\ell,i,j}$ such that
\[x^\ell=\sum_{i_1,i_2:i_1q+i_2\leq\ell,i_2<q}\beta_{\ell,i_1,i_2}g(x)^{i_1}x^{i_2}.\]
By comparing coefficients of $x^\ell$ we see that $\beta_{\ell,\lfloor\ell/q\rfloor,\ell\bmod q}=1$.
\paragraph{Claim 44.} For every table $T:\mathbb F_q^m\to\Sigma_{m,s}$ there exists a degree $q(s-1)+(q-1)m$ polynomial such that $\mathrm{EVAL}(P)=T$.
\paragraph{Proof.} By Claim 41, $T$ may be written as $\mathrm{EVAL}(P)$, where $P$ is an $\mathbb F_q$-combination of basis elements in $B_{m,s}$. By the definition of $B_{m,s}$, any basis element is of the form $g(\boldsymbol X)^{\boldsymbol I}\cdot\prod_{k=1}^mX_k^{j_k}$ where $w(\boldsymbol I)<s$ and $0\leq j_k<q$. As such, $\deg(P)\leq\deg(g(\boldsymbol X)^{\boldsymbol I})+\deg(\prod_{k=1}^mX_k^{j_k})\leq q(s-1)+(q-1)m$.\hfill$\triangleleft$
We will later need:
\paragraph{Lemma 45.} Let $A\in\mathbb F_q[t]$ and $B\in\mathbb F_q[r]$. Then,
\[\deg_t(A(t)B(r)\bmod I_{2,s})\leq\deg_t(A).\]
\paragraph{Proof.} Express $A$ in the basis $\{g(t)^it^j\}_{i\in\mathbb N,j<q}$,
\[A(t)=\sum_{(i,j)\in\mathcal A}\alpha_{i,j}g(t)^it^j,\]
and similarly for $B$,
\[B(r)=\sum_{(k,\ell)\in\mathcal B}\beta_{k,\ell}g(r)^kr^\ell,\]
for some sets $\mathcal A,\mathcal B$, $\alpha_{i,j},\beta_{k,\ell}\in\mathbb F_q$. Then,
\[A(t)B(r)\bmod I_{2,s}=\sum_{(i,j)\in\mathcal A,(k,\ell)\in\mathcal B:i+k<s}\alpha_{i,j}\beta_{k,\ell}g(t)^ig(r)^kt^jr^\ell.\]
Suppose $\deg_t(A(t)B(r)\bmod I_{2,s})$ is achieved on the monomial of $(i,j)\in\mathcal A,(k,\ell)\in\mathcal B$. In particular, this degree in $t$ is already achieved in $A$ for the monomial $(i,j)$. Thus, $\deg_t(AB)\leq\deg_t(A)$.\hfill$\triangleleft$
Next, we would like to give a purely algebraic criteria, which states when exactly a table belongs to the code $\mathrm{MRM}(q,m,d,s)$.
\paragraph{Definition 46.} Let $T\in\Sigma_{m,s}^{q^m}$ be a table. By Corollary 40 and Claim 41 there is a unique element $P_T\in\mathbb F[\boldsymbol X]/I_{m,s}$ such that $\mathrm{EVAL}_{m,s}(P_T)=T$. We call $P_T$ the representing polynomial of $T$.
\paragraph{Lemma 47.} Assume $d<sq$. Let $T\in\Sigma_{m,s}^{q^m}$ be a table, and $P_T$ its representing polynomial. Then
\[T\in\mathrm{MRM}(q,m,d,s)\Longleftrightarrow\deg(P_T)\leq d.\]
\paragraph{Proof.} First, assume $\deg(P_T)\leq d$. Then, by definition, since $\mathrm{EVAL}_{m,s}(P_T)=T$, we have $T\in\mathrm{MRM}(q,m,d,s)$. For the other direction, assume $T\in\mathrm{MRM}(q,m,d,s)$. Then there is some $Q\in\mathbb F[\boldsymbol X]$ of total degree $\leq d$ such that $\mathrm{EVAL}_{m,s}(Q)=T$. As $\deg(Q)\leq d<sq$, $Q$ is $I_{m,s}$ reduced and by Corollary 40 $Q$ is the representing polynomial of $T$.\hfill$\triangleleft$
\section*{Imported RM tester and numerical transfer (author context)}
\paragraph{Theorem 2 ([12]).} Let $k=\lceil(d+1)/(q-q/p)\rceil$. Given $f:\mathbb F_q^m\to\mathbb F_q$ let:
\begin{itemize}
\item $\mathrm{REJ}_{k,d}(f)$ denotes the rejection probability of the test, namely, the probability over a random $k$ dimensional affine space, that $f$ restricted to the space cannot be explained by a degree $k$ polynomial, and,
\item $\delta(f,\mathrm{RM}(q,m,d))$ be the distance of $f$ from the code $\mathrm{RM}(q,m,d)$.
\end{itemize}
Then
\[\mathrm{REJ}_{k,d}(f)\geq\min\{\alpha_0\cdot\delta(f,\mathrm{RM}(q,m,d)),c_0\}.\]
where $\alpha_0=q^k/2$ and $c_0=1/(2(k+1)q^{k+1})$.
Combining the result with Theorem 2 we get:
\paragraph{Corollary 52.} Let $d,q,s<q$ be positive integers and let $t=t_{q,d}=\lceil(d+1)/(q-q/p)\rceil$ be the RM testing dimension. Then,
\[\mathrm{REJ}_{t,d}^{\mathrm{MRM}}(T)\geq\min\left\{\frac13\left(1-\frac{s-1}q\right)\delta(T,\mathrm{MRM}(q,m,d,s)),\frac1{2(t+1)q^{t+1}}\right\}.\]
\paragraph{Proof.} Theorem 2 from [12] ensures that
\[\mathrm{REJ}_{t,d}(f)\geq\min\left\{\frac{q^t}2\delta(f,\mathrm{RM}(q,m,d)),\frac1{2(t+1)q^{t+1}}\right\}.\]
Using Theorem 51 we see that
\[\alpha'\geq\left(1-\frac{s-1}q\right)\frac1{1+q^{d/(q-1)}\frac1\alpha}\geq\frac13\left(1-\frac{s-1}q\right).\]
using the fact that $d/(q-1)\leq t$ and $q^{d/(q-1)}/\alpha\leq2$.\hfill$\triangleleft$
In particular, for $d<(q-q/p)$ the line test is a local characterization, as $t_{q,d}=1$.
We note that the assumption $s\leq d$ is quite natural, as derivatives with order higher than the degree must be identically zero. In contrast, it is not clear whether the assumption $s<q$ is indeed required, and we leave it for future study.

% Faithful official CCC2022 article14 Sections6–7.
\section*{6 The plane test is a characterization}
In this section we show that the multiplicity code $\mathrm{MRM}(q,m,d,s)$ can be characterized by restrictions to planes. Let $\mathbb F_q$ be a field and let $m,s\leq d$ be positive integers. For $\boldsymbol a,\boldsymbol b,\boldsymbol c\in\mathbb F_q^m$\footnote{In this section, we omit the requirement that $\boldsymbol a,\boldsymbol b$ be independent. If some degenerate plane shows that a table is not a low degree polynomial, then some actual plane will too.} define: $\ell_{\boldsymbol a,\boldsymbol b,\boldsymbol c}:\mathbb F_q^2\to\mathbb F_q^m$ by:
\[\ell_{\boldsymbol a,\boldsymbol b,\boldsymbol c}(t,r)=\boldsymbol at+\boldsymbol br+\boldsymbol c.\]
From Lemma 18 and Definition 48 we see that:
\paragraph{Theorem 59 (Completeness).} Suppose $d<sq-1$. If a table $T\in\Sigma_{m,s}^{q^m}$ satisfies $T\in\mathrm{MRM}(q,m,d,s)$ then for all $\boldsymbol a,\boldsymbol b,\boldsymbol c\in\mathbb F_q^m$,
\[\phi_{(\boldsymbol a,\boldsymbol b)}\circ T\circ\ell_{\boldsymbol a,\boldsymbol b,\boldsymbol c}\in\mathrm{MRM}(q,2,d,s).\]
The main challenge is proving the converse:
\paragraph{Theorem 60 (Soundness).} Suppose $q$ is a power of the prime $p$, $q\geq s$ and $d<q(s-1/p)$. If a table $T\in\Sigma_{m,2}^{q^m}$ satisfies that for all $\boldsymbol a,\boldsymbol b,\boldsymbol c\in\mathbb F_q^m$, $\phi_{(\boldsymbol a,\boldsymbol b)}\circ T\circ\ell_{\boldsymbol a,\boldsymbol b,\boldsymbol c}\in\mathrm{MRM}(q,2,d,s)$ then $T\in\mathrm{MRM}(q,m,d,s)$.
We define the vector space of tables which pass the test:
\[V_{m,d,s}=\{T\in\Sigma_{m,s}^{q^m}\mid\forall\boldsymbol a,\boldsymbol b,\boldsymbol c\in\mathbb F_q^m\ \phi_{(\boldsymbol a,\boldsymbol b)}\circ T\circ\ell_{\boldsymbol a,\boldsymbol b,\boldsymbol c}\in\mathrm{MRM}(q,2,d,s)\}.\tag{12}\]
We denote by $C_{m,d,s}=V_{m,d,s}\setminus\mathrm{MRM}_{m,d,s}$, the set of tables which cheat the test. We would like to show that $C_{m,d,s}=\emptyset$. Assume towards a contradiction that there is a table $T\in C_{m,d,s}$. By Claim 41, $T$ can be realized (uniquely) as an element $P$ of the quotient space $\mathbb F_q[\boldsymbol X]/I_{m,s}$. We use the basis $B_{m,s}$ from Equation (5) to write $P$ in the form
\[P(\boldsymbol X)=\sum_{(\boldsymbol I,\boldsymbol J)\in M_{s,q}}\alpha_{\boldsymbol I,\boldsymbol J}\cdot g(\boldsymbol X)^{\boldsymbol I}\boldsymbol X^{\boldsymbol J},\]
where $\alpha_{\boldsymbol I,\boldsymbol J}\in\mathbb F_q$ and $(\boldsymbol I,\boldsymbol J)\in M_{s,q}$ iff $w(\boldsymbol I)<s$ and $J_k<q$ for every $1\leq k\leq m$. Since $T\notin\mathrm{MRM}_{m,d,s}$ we have $\deg(P)>d$. This means there must be some $\boldsymbol I$ and $\boldsymbol J$ such that $\alpha_{\boldsymbol I,\boldsymbol J}\neq0$ and
\[w(\boldsymbol I)q+w(\boldsymbol J)>d.\tag{13}\]
We may assume that every $\boldsymbol I,\boldsymbol J$ for which $\alpha_{\boldsymbol I,\boldsymbol J}\neq0$ satisfy Equation (13). This is since the test is linear, and any degree $\leq d$ terms have no effect on whether $P$ passes the test or not.
We use the following monomial order $\succ_w$ on $B_{m,s}$:
\begin{enumerate}
\item First order monomials according to $w(\boldsymbol I)+w(\boldsymbol J)$,
\item Then order monomials according to $w(\boldsymbol I)$,
\item Finally, order monomials according to the lexicographic order on $\boldsymbol I,\boldsymbol J$.
\end{enumerate}
We emphasize that $\succ_w$ is not a monomial order in the sense of Grobner bases, and we make no use of it in that sense.
Let $(\boldsymbol I_{\max},\boldsymbol J_{\max})$ be s.t. $g(\boldsymbol X)^{\boldsymbol I_{\max}}\cdot\boldsymbol X^{\boldsymbol J_{\max}}$ is a maximal monomial of $P$ according to $\succ_w$.
For $\boldsymbol a,\boldsymbol b,\boldsymbol c\in\mathbb F^m$ the restriction of $P$ to the plane defined by $\boldsymbol a,\boldsymbol b,\boldsymbol c$ is $P|_{P_{\boldsymbol a,\boldsymbol b,\boldsymbol c}}(t,r)=P(\boldsymbol at+\boldsymbol br+\boldsymbol c)$. Expressing $P|_{P_{\boldsymbol a,\boldsymbol b,\boldsymbol c}}(t,r)$ in the basis $B_{2,\infty}$:
\begin{align*}
P|_{P_{\boldsymbol a,\boldsymbol b,\boldsymbol c}}(t,r)&=\sum_{i\in\mathbb N,j\in\mathbb N,k,\ell<q}A_{i,j,k,\ell}(\boldsymbol a,\boldsymbol b,\boldsymbol c)\cdot g(t)^ig(r)^jt^kr^\ell,\\
P|_{P_{\boldsymbol a,\boldsymbol b,\boldsymbol c}}(t,r)\bmod I_{2,s}&=\sum_{i+j<s,k,\ell<q}A_{i,j,k,\ell}(\boldsymbol a,\boldsymbol b,\boldsymbol c)\cdot g(t)^ig(r)^jt^kr^\ell.
\end{align*}
We view $A_{i,j,k,\ell}(\boldsymbol a,\boldsymbol b,\boldsymbol c)$ as a polynomial in the variables $\boldsymbol a,\boldsymbol b,\boldsymbol c$.
\paragraph{Lemma 61.} For every partition $\boldsymbol J_{\max}=\boldsymbol J_{\max}^b+\boldsymbol J_{\max}^c$ such that:
\begin{itemize}
\item $q\cdot w(\boldsymbol I_{\max})+w(\boldsymbol J_{\max}^b)\leq qs-1$, and,
\item $\binom{\boldsymbol J_{\max}}{\boldsymbol J_{\max}^b}\neq0\bmod p$,
\end{itemize}
the monomial $\boldsymbol a^{\boldsymbol I_{\max}}\boldsymbol b^{\boldsymbol J_{\max}^b}\boldsymbol c^{\boldsymbol J_{\max}^c}$ appears with a non-zero coefficient at
\[A_{w(\boldsymbol I_{\max}),\lfloor w(\boldsymbol J_{\max}^b)/q\rfloor,0,w(\boldsymbol J_{\max}^b)\bmod q}(\boldsymbol a,\boldsymbol b,\boldsymbol c)\bmod I_{3m,1}.\]
Where the ideal $I_{3m,1}$ above is in the variables $a_1,\ldots,a_m,b_1,\ldots,b_m,c_1,\ldots,c_m$.
\paragraph{Proof.} We expand $P|_{P_{\boldsymbol a,\boldsymbol b,\boldsymbol c}}(t,r)$. First, $P(\boldsymbol X)=\sum_{(\boldsymbol I,\boldsymbol J)\in M_{s,q}}\alpha_{\boldsymbol I,\boldsymbol J}\cdot g(\boldsymbol X)^{\boldsymbol I}\boldsymbol X^{\boldsymbol J}$. Thus
\begin{align*}
P|_{P_{\boldsymbol a,\boldsymbol b,\boldsymbol c}}(t,r)
&=\sum_{(\boldsymbol I,\boldsymbol J)\in M_{s,q}}\alpha_{\boldsymbol I,\boldsymbol J}\cdot g(\boldsymbol at+\boldsymbol br+\boldsymbol c)^{\boldsymbol I}(\boldsymbol at+\boldsymbol br+\boldsymbol c)^{\boldsymbol J}\\
&=\sum_{(\boldsymbol I,\boldsymbol J)\in M_{s,q}}\alpha_{\boldsymbol I,\boldsymbol J}\prod_{k=1}^m(g(a_kt+b_kr+c_k))^{I_k}\cdot\prod_{i=k}^m(a_kt+b_kr+c_k)^{J_k}\\
&=\sum_{(\boldsymbol I,\boldsymbol J)\in M_{s,q}}\alpha_{\boldsymbol I,\boldsymbol J}\prod_{k=1}^m(g(t)a_k+g(r)b_k)^{I_k}\cdot\prod_{i=k}^m(a_kt+b_kr+c_k)^{J_k}\\
&=\sum_{(\boldsymbol I,\boldsymbol J)\in M_{s,q}}\alpha_{\boldsymbol I,\boldsymbol J}\sum_{\substack{\boldsymbol I_a+\boldsymbol I_b=\boldsymbol I\\\boldsymbol J_a+\boldsymbol J_b+\boldsymbol J_c=\boldsymbol J}}\binom{\boldsymbol I}{\boldsymbol I_a}\binom{\boldsymbol J}{\boldsymbol J_a,\boldsymbol J_b,\boldsymbol J_c}\\
&\qquad\boldsymbol a^{\boldsymbol I_a+\boldsymbol J_a}\boldsymbol b^{\boldsymbol I_b+\boldsymbol J_b}\boldsymbol c^{\boldsymbol J_c}\cdot g(t)^{w(\boldsymbol I_a)}g(r)^{w(\boldsymbol I_b)}t^{w(\boldsymbol J_a)}r^{w(\boldsymbol J_b)}.\tag{14}
\end{align*}
We expand $t^{w(\boldsymbol J_a)}$ and $r^{w(\boldsymbol J_b)}$ in the basis $B_{1,s}$ as in Example 43 to get:
\begin{align*}
P|_{P_{\boldsymbol a,\boldsymbol b,\boldsymbol c}}(t,r)
&=\sum_{(\boldsymbol I,\boldsymbol J)\in M_{s,q}}\alpha_{\boldsymbol I,\boldsymbol J}\cdot\sum_{\substack{\boldsymbol I_a+\boldsymbol I_b=\boldsymbol I\\\boldsymbol J_a+\boldsymbol J_b+\boldsymbol J_c=\boldsymbol J}}\binom{\boldsymbol I}{\boldsymbol I_a}\binom{\boldsymbol J}{\boldsymbol J_a\ \boldsymbol J_b\ \boldsymbol J_c}\\
&\quad\sum_{\substack{i_1,i_2\\i_1q+i_2\leq w(\boldsymbol J_a)\\i_2<q}}\sum_{\substack{j_1,j_2\\j_1q+j_2\leq w(\boldsymbol J_b)\\j_2<q}}\beta_{w(\boldsymbol J_a),i_1,i_2}\beta_{w(\boldsymbol J_b),j_1,j_2}\\
&\quad\boldsymbol a^{\boldsymbol I_a+\boldsymbol J_a}\boldsymbol b^{\boldsymbol I_b+\boldsymbol J_b}\boldsymbol c^{\boldsymbol J_c}\cdot g(t)^{w(\boldsymbol I_a)+i_1}t^{i_2}\cdot g(r)^{w(\boldsymbol I_b)+j_1}r^{j_2}.\tag{15}
\end{align*}
We now have a representation in the basis $B_{2,\infty}$.
Set:
\begin{align*}
\Delta_{rq}&=\left\lfloor\frac{w(\boldsymbol J_{\max}^b)}q\right\rfloor,\\
\Delta_t&=0,\\
\Delta_r&=w(\boldsymbol J_{\max}^b)\bmod q.\tag{16}
\end{align*}
We wish to see for which choice of values $\boldsymbol I,\boldsymbol J,\boldsymbol I_a,\boldsymbol I_b,\boldsymbol J_a,\boldsymbol J_b,\boldsymbol J_c,i_1,i_2,j_1,j_2$ in Equation (14), $\boldsymbol a^{\boldsymbol I_{\max}}\boldsymbol b^{\boldsymbol J_{\max}^b}\boldsymbol c^{\boldsymbol J_{\max}^c}$ appears as a coefficient of $A_{(w(\boldsymbol I_{\max}),\Delta_{rq},\Delta_t,\Delta_r)}\bmod I_{3m,1}$. We must have
\begin{align*}
\boldsymbol I_{\max}&=_{\mathbb F_q}\boldsymbol I_a+\boldsymbol J_a&&\text{By comparing the powers of }\boldsymbol a,\text{ remembering }\bmod I_{3m,1},\\
\boldsymbol J_{\max}^b&=_{\mathbb F_q}\boldsymbol I_b+\boldsymbol J_b&&\text{By comparing the powers of }\boldsymbol b,\text{ remembering }\bmod I_{3m,1},\\
\boldsymbol J_{\max}^c&=_{\mathbb F_q}\boldsymbol J_c&&\text{By comparing the powers of }\boldsymbol c,\text{ remembering }\bmod I_{3m,1},
\end{align*}
where $=_{\mathbb F_q}$ was defined in Definition 36.
By Claim 38 together with $s<q$ we have
\begin{align*}
w(\boldsymbol I_{\max})&\leq w(\boldsymbol I_a)+w(\boldsymbol J_a),\tag{17}\\
w(\boldsymbol J_{\max}^b)&\leq w(\boldsymbol I_b)+w(\boldsymbol J_b),\tag{18}\\
\boldsymbol J_{\max}^c&=\boldsymbol J_c.\tag{19}
\end{align*}
It follows that
\begin{align*}
w(\boldsymbol I_{\max})+w(\boldsymbol J_{\max})
&=w(\boldsymbol I_{\max})+w(\boldsymbol J_{\max}^b)+w(\boldsymbol J_{\max}^c)\\
&\leq w(\boldsymbol I_a)+w(\boldsymbol J_a)+w(\boldsymbol I_b)+w(\boldsymbol J_b)+w(\boldsymbol J_c)\\
&=w(\boldsymbol I)+w(\boldsymbol J).
\end{align*}
As $(\boldsymbol I_{\max},\boldsymbol J_{\max})$ is maximal it follows that
\[w(\boldsymbol I_{\max})+w(\boldsymbol J_{\max})=w(\boldsymbol I)+w(\boldsymbol J).\tag{20}\]
This, in turn, implies that both inequalities in Equations (17) and (18) are in fact equalities, i.e.,
\begin{align*}
w(\boldsymbol I_{\max})&=w(\boldsymbol I_a)+w(\boldsymbol J_a),\tag{21}\\
w(\boldsymbol J_{\max}^b)&=w(\boldsymbol I_b)+w(\boldsymbol J_b).\tag{22}
\end{align*}
We now look at
\[\deg_t\stackrel{\mathrm{def}}=\deg_t(g(t\boldsymbol a)^{\boldsymbol I_a}\cdot(t\boldsymbol a)^{\boldsymbol J_a}\cdot g(r\boldsymbol b)^{\boldsymbol I_b}\cdot(r\boldsymbol b)^{\boldsymbol J_b}\cdot\boldsymbol c^{\boldsymbol J_c}\bmod I_{2,s}).\]
On the one hand, we look for the monomial $\boldsymbol a^{\boldsymbol I_{\max}}\boldsymbol b^{\boldsymbol J_{\max}^b}\boldsymbol c^{\boldsymbol J_{\max}^c}$ in $A_{w(\boldsymbol I_{\max}),\Delta_{rq},\Delta_t,\Delta_r}\bmod I_{3m,1}$, and so we should have $\deg_t=q\cdot w(\boldsymbol I_{\max})+\Delta_t$. On the other hand, by Lemma 45,
\begin{align*}
\deg_t&=\deg_t(g(t\boldsymbol a)^{\boldsymbol I_a}(t\boldsymbol a)^{\boldsymbol J_a}g(r\boldsymbol b)^{\boldsymbol I_b}(r\boldsymbol b)^{\boldsymbol J_b}\boldsymbol c^{\boldsymbol J_c}\bmod I_{2,s})\\
&\leq\deg_t(g(t\boldsymbol a)^{\boldsymbol I_a}(t\boldsymbol a)^{\boldsymbol J_a})\\
&=q\cdot w(\boldsymbol I_a)+w(\boldsymbol J_a).
\end{align*}
Thus,
\[q\cdot w(\boldsymbol I_{\max})+\Delta_t\leq q\cdot w(\boldsymbol I_a)+w(\boldsymbol J_a).\]
Together with Equation (21) we see that $q\cdot w(\boldsymbol I_a)+q\cdot w(\boldsymbol J_a)+\Delta_t\leq q\cdot w(\boldsymbol I_a)+w(\boldsymbol J_a)$, and therefore $q\cdot w(\boldsymbol J_a)+\Delta_t\leq w(\boldsymbol J_a)$ which is possible iff $w(\boldsymbol J_a)=0$ (and $\Delta_t=0$, which is indeed true, see Equation (16)). Now, $w(\boldsymbol J_a)=0$ implies:
\[w(\boldsymbol I_a)=w(\boldsymbol I_{\max}),\text{ and,}\qquad\boldsymbol J_a=\emptyset.\]
We saw in Equation (20) that $w(\boldsymbol I)+w(\boldsymbol J)=w(\boldsymbol I_{\max})+w(\boldsymbol J_{\max})$. If $\boldsymbol I_b\neq\emptyset$ then
\[w(\boldsymbol I)=w(\boldsymbol I_a)+w(\boldsymbol I_b)>w(\boldsymbol I_a)=w(\boldsymbol I_a)+w(\boldsymbol J_a)=w(\boldsymbol I_{\max}).\]
Thus, $(\boldsymbol I,\boldsymbol J)\succ_w(\boldsymbol I_{\max},\boldsymbol J_{\max})$ in contradiction to the maximality of $(\boldsymbol I_{\max},\boldsymbol J_{\max})$. We conclude that $\boldsymbol I_b=\emptyset$.
As $(\boldsymbol I_{\max},\boldsymbol J_{\max})\in M_{s,q}$ we have $(\boldsymbol J_{\max})_k\leq q-1$ for every $k\in[m]$. Thus, $\boldsymbol J_{\max}^b=_{\mathbb F_q}\boldsymbol J_b$, $w(\boldsymbol J_{\max}^b)=w(\boldsymbol J_b)$ and $\boldsymbol J_{\max}^b$ is already reduced. Together this implies that $\boldsymbol J_{\max}^b=\boldsymbol J_b$.
Finally, we use the hypothesis that $q\geq s$. We have, $(\boldsymbol I_{\max})_k<s\leq q$ for all $k\in[m]$. Thus, $\boldsymbol I_{\max}=_{\mathbb F_q}\boldsymbol I_a$, $w(\boldsymbol I_{\max})=w(\boldsymbol I_a)$ and $\boldsymbol I_{\max}$ is already reduced. Together this implies that that $\boldsymbol I_{\max}=\boldsymbol I_a$.
Thus,
\[\boldsymbol I_{\max}=\boldsymbol I_a=\boldsymbol I_a+\boldsymbol I_b=\boldsymbol I,\qquad\boldsymbol J_{\max}^b=\boldsymbol I_b+\boldsymbol J_b=\boldsymbol J_b.\]
Altogether,the only term that may possibly contribute $\boldsymbol a^{\boldsymbol I_{\max}}\boldsymbol b^{\boldsymbol J_{\max}^b}\boldsymbol c^{\boldsymbol J_{\max}^c}$ to $A_{w(\boldsymbol I_{\max}),\Delta_{rq},\Delta_t,\Delta_r}\bmod I_{3m,1}$ is
\[(\boldsymbol I_a,\boldsymbol I_b,\boldsymbol J_a,\boldsymbol J_b,\boldsymbol J_c)=(\boldsymbol I_{\max},\emptyset,\emptyset,\boldsymbol J_{\max}^b,\boldsymbol J_{\max}^c).\]
Also, the tuple $(\boldsymbol I_a,\boldsymbol I_b,\boldsymbol J_a,\boldsymbol J_b,\boldsymbol J_c)=(\boldsymbol I_{\max},\emptyset,\emptyset,\boldsymbol J_{\max}^b,\boldsymbol J_{\max}^c)$ contributes
\begin{multline*}
\alpha_{\boldsymbol I_{\max},\boldsymbol J_{\max}}\cdot\binom{\boldsymbol J_{\max}}{\boldsymbol J_{\max}^b}\sum_{\substack{j_1,j_2\\j_1q+j_2\leq w(\boldsymbol J_{\max}^b)\\j_2<q}}\beta_{w(\boldsymbol J_{\max}^b),j_1,j_2}\\
\boldsymbol a^{\boldsymbol I_{\max}}\boldsymbol b^{\boldsymbol J_{\max}^b}\boldsymbol c^{\boldsymbol J_{\max}^c}\cdot g(t)^{w(\boldsymbol I_{\max})}\cdot g(r)^{j_1}r^{j_2}
\end{multline*}
to the term in Equation (14).
Notice that $w(\boldsymbol I_{\max})+j_1<s$, for otherwise $q\cdot w(\boldsymbol I_{\max})+w(\boldsymbol J_{\max}^b)\geq q(w(\boldsymbol I_{\max})+j_1)\geq qs$ in contradiction to the hypothesis. Thus the term is already $I_{2,s}$ reduced. The contribution to $A_{w(\boldsymbol I_{\max}),\Delta_{rq},0,\Delta_r}\bmod I_{3m,1}$ occurs for $(j_1,j_2)$ such that $j_1=\Delta_{rq}$ and $j_2=\Delta_r$. Thus, in the sum in Equation (14) there is exactly one possible way to contribute $\boldsymbol a^{\boldsymbol I_{\max}}\boldsymbol b^{\boldsymbol J_{\max}^b}\boldsymbol c^{\boldsymbol J_{\max}^c}$ to $A_{w(\boldsymbol I_{\max}),\Delta_{rq},0,\Delta_r}\bmod I_{3m,1}$, and this is when
\begin{align*}
(\boldsymbol I_a,\boldsymbol I_b,\boldsymbol J_a,\boldsymbol J_b,\boldsymbol J_c)&=(\boldsymbol I_{\max},\emptyset,\emptyset,\boldsymbol J_{\max}^b,\boldsymbol J_{\max}^c),\text{ and}\\
(i_1,i_2,j_1,j_2)&=(0,0,\Delta_{rq},\Delta_r).
\end{align*}
The coefficient of this term is:
\[\alpha_{\boldsymbol I_{\max},\boldsymbol J_{\max}}\cdot\binom{\boldsymbol J_{\max}}{\boldsymbol J_{\max}^b}\cdot\beta_{w(\boldsymbol J_{\max}^b),\Delta_{rq},\Delta_r}.\]
$\alpha_{\boldsymbol I_{\max},\boldsymbol J_{\max}}\neq0$. By assumption $\binom{\boldsymbol J_{\max}}{\boldsymbol J_{\max}^b}$ is non-zero. $\beta_{\ell,\lfloor\ell/q\rfloor,\ell\bmod q}=1$ (see Example 43) and taking $\ell=w(\boldsymbol J_{\max}^b)$ shows the coefficient is non-zero. As there is a unique term contributing the monomial with a non-zero coefficient, the monomial cannot cancel in $A_{w(\boldsymbol I_{\max}),\Delta_{rq},\Delta_t,\Delta_r}\bmod I_{3m,1}$ and $A_{w(\boldsymbol I_{\max}),\Delta_{rq},\Delta_t,\Delta_r}\bmod I_{3m,1}$ is non-zero.\hfill$\triangleleft$
\paragraph{Lemma 62.} There is a partition $\boldsymbol J_{\max}=\boldsymbol J_{\max}^b+\boldsymbol J_{\max}^c$, such that
\begin{itemize}
\item $d<q\cdot w(\boldsymbol I_{\max})+w(\boldsymbol J_{\max}^b)\leq qs-1$, and,
\item $\binom{\boldsymbol J_{\max}}{\boldsymbol J_{\max}^b}\bmod p\neq0$.
\end{itemize}
\paragraph{Proof.} Suppose $q=p^w$ where $p$ is prime (if $q$ is prime then $p=q$ and $w=1$). We choose $\boldsymbol J_{\max}^b$ as follows. We go over $k=1,\ldots,m$ and find the first $k_0\geq0$ such that $q\cdot w(\boldsymbol I_{\max})+\sum_{k=0}^{k_0}(\boldsymbol J_{\max})_k>d$. There must be some $k_0\leq m$ like that since $qw(\boldsymbol I_{\max})+w(\boldsymbol J_{\max})>d$. We set:
\begin{itemize}
\item $(\boldsymbol J_{\max}^b)_k=(\boldsymbol J_{\max})_k$, for $k=1,\ldots,k_0-1$, and,
\item $(\boldsymbol J_{\max}^b)_k=0$, for $k=k_0+1,\ldots,m$.
\end{itemize}
Set $v=(\boldsymbol J_{\max})_{k_0}$. There is a first value $0<v'\leq v$ such that
\[qw(\boldsymbol I_{\max})+\sum_{k=1}^{k_0-1}(\boldsymbol J_{\max})_k+v'>d.\]
We express $v=(\boldsymbol J_{\max})_{k_0}$ in base $p$: $v=\sum_{\ell=0}^{w-1}v_\ell\cdot p^\ell$. We let $v''$ be the first integer such that:
\begin{itemize}
\item $v''\geq v'$, and,
\item If we express $v''$ in base $p$ as $v''=\sum_{\ell=0}^{w-1}v''_\ell\cdot p^\ell$ then $v''_\ell\leq v_\ell$ for every $\ell$.
\end{itemize}
\paragraph{Claim 63.} $v'\leq v''\leq\min\{v,v'+p^{w-1}-1\}$.
\paragraph{Proof.} Notice that $v$ respects the conditions that we need, and so if $v\leq v'+p^{w-1}-1$ the claim holds. Otherwise, $v\geq v'+p^{w-1}-1$. Then, $v''\leq v'+p^{w-1}-1$ because we can increase $v'$ by setting all the lower bits in the $p$-representation to $0$, while increasing the most significant bit (that is multiplied by $p^{w-1}$) by $1$.\hfill$\triangleleft$
Thus, by Lucas theorem (see Equation (4))
\[\binom v{v''}\bmod p=\prod_{\ell=0}^{w-1}\binom{v_\ell}{v''_\ell}\neq0.\]
Having that, we let $(\boldsymbol J_{\max}^b)_{k_0}=v''$. We have, $\boldsymbol J_{\max}^b\leq\boldsymbol J_{\max}$, $qw(\boldsymbol I_{\max})+w(\boldsymbol J_{\max}^b)>d$ and $\binom{\boldsymbol J_{\max}}{\boldsymbol J_{\max}^b}\bmod p\neq0$. Also,
\[qw(\boldsymbol I_{\max})+w(\boldsymbol J_{\max}^b)\leq d+p^{w-1}\leq qs-1,\]
because $d\leq qs-p^{w-1}-1=q(s-1/p)-1$, completing the proof of the lemma.\hfill$\triangleleft$
We are now ready to prove Theorem 60.
\paragraph{Proof.} Fix a partition $\boldsymbol J_{\max}=\boldsymbol J_{\max}^b+\boldsymbol J_{\max}^c$ as in Lemma 62. Let
\begin{align*}
\deg_r&=\deg_r(g(t)^{w(\boldsymbol I_{\max})}r^{w(\boldsymbol J_{\max}^b)}\bmod I_{2,s})\\
&=\deg_r(g(t)^{w(\boldsymbol I_{\max})}r^{w(\boldsymbol J_{\max}^b)})=w(\boldsymbol J_{\max}^b).
\end{align*}
Define $\Delta_{rq}=\lfloor w(\boldsymbol J_{\max}^b)/q\rfloor$, $\Delta_t=0$ and $\Delta_r=w(\boldsymbol J_{\max}^b)\bmod q$. By Lemma 61 we know that
\[A_{w(\boldsymbol I_{\max}),\Delta_{rq},\Delta_t,\Delta_r}(\boldsymbol a,\boldsymbol b,\boldsymbol c)\bmod I_{3m,1}\neq0.\]
Thus, there exist $\boldsymbol a_0,\boldsymbol b_0,\boldsymbol c_0$ such that
\[A_{w(\boldsymbol I_{\max}),\Delta_{rq},\Delta_t,\Delta_r}(\boldsymbol a_0,\boldsymbol b_0,\boldsymbol c_0)\neq0.\]
We look at the test $\boldsymbol a_0,\boldsymbol b_0,\boldsymbol c_0$. We have
\[P|_{P_{\boldsymbol a,\boldsymbol b,\boldsymbol c}}(r,t)\bmod I_{2,s}=\sum_{i+j<s,k,\ell<q}A_{i,j,k,\ell}(\boldsymbol a,\boldsymbol b,\boldsymbol c)g(t)^ig(r)^jt^kr^\ell.\]
As $A_{w(\boldsymbol I_{\max}),\Delta_{rq},\Delta_t,\Delta_r}(\boldsymbol a_0,\boldsymbol b_0,\boldsymbol c_0)\neq0$ we see that
\begin{align*}
\deg(P|_{P_{\boldsymbol a_0,\boldsymbol b_0,\boldsymbol c_0}}\bmod I_{2,s})
&\geq q\cdot w(\boldsymbol I_{\max})+q\cdot\Delta_{rq}+\Delta_t+\Delta_r\\
&=q\cdot w(\boldsymbol I_{\max})+w(\boldsymbol J_{\max}^b)>d.
\end{align*}
Thus, by Lemma 47, the test $(\boldsymbol a_0,\boldsymbol b_0,\boldsymbol c_0)$ rejects.\hfill$\triangleleft$
\section*{7 Planes give a local MRM tester}
We restate Theorem 9:

\paragraph{Theorem 64.} Suppose $q$ is a prime power, $s\leq q$ and $d<q(s-1/p)$. Let $T:\mathbb F_q^m\to\Sigma_{m,s}$ be a table and let $\delta=\delta(T,\mathrm{MRM}(q,m,d,s))$. Then
\[\mathrm{REJ}_{2,d}^{\mathrm{MRM}}(T)\geq\min\{\alpha\delta,c\}\]
with $\alpha=\Omega(q^{-6s+5})$ and $c=\Omega(q^{-8s+4})$.

\paragraph{Proof.} We remind the reader that
\[H_k=\{\boldsymbol h=(\boldsymbol h_0,\boldsymbol h_1,\ldots,\boldsymbol h_k)\mid\boldsymbol h_0,\ldots,\boldsymbol h_k\in\mathbb F_q^m,\dim(\operatorname{span}\{\boldsymbol h_1,\ldots,\boldsymbol h_k\})=k\}.\]
Let us say that $\boldsymbol v\in H_k$ is bad, if the $k$-dimensional test with $\boldsymbol v$ rejects, i.e., $(\phi_{\boldsymbol v}\circ T)|_{A_{\boldsymbol v}}\notin\mathrm{MRM}(q,k,d,s)$. We let $t=\lceil(d+1)/(q-q/p)\rceil$ be the $\mathrm{RM}(q,m,d,s)$ testing dimension. Selecting a uniform $\boldsymbol u\in H_2$ is the same as selecting $\boldsymbol h\sim H_t$ and then a uniform $\boldsymbol u\in H_2$ such that $A_{\boldsymbol u}\subset A_{\boldsymbol h}$. Thus,
\begin{align*}
\mathrm{REJ}_{2,d}^{\mathrm{MRM}}(T)&=\Pr_{\boldsymbol u\sim H_2}(\boldsymbol u\text{ is bad})\\
&\geq\Pr_{\boldsymbol h\in H_t}(\boldsymbol h\text{ is bad})\cdot\Pr_{\boldsymbol u\in H_2:A_{\boldsymbol u}\subset A_{\boldsymbol h}}(\boldsymbol u\text{ is bad}\mid\boldsymbol h\text{ is bad}).
\end{align*}
\begin{itemize}
\item By Corollary 52, the probability of picking a bad $\boldsymbol h\in H_t$ is:
\[\mathrm{REJ}_{t,d}^{\mathrm{MRM}}(T)\geq\min\left\{\frac\delta{3q},\frac1{2(t+1)q^{t+1}}\right\}.\]
\item By Theorem 60 we know that for any bad $\boldsymbol h\in H_t$ there is at least one $\boldsymbol u\in H_2$ such that $A_{\boldsymbol u}$ is contained $A_{\boldsymbol h}$ and $\boldsymbol u$ is bad. Furthermore, if $A_{\boldsymbol u}=A_{\boldsymbol u'}$ then $\boldsymbol u$ is bad iff $\boldsymbol u'$ is bad. We look at $A_{\boldsymbol u}$. There are $(q^2-1)(q^2-q)q^2$ different $\boldsymbol u'\in H_2$ such that $A_{\boldsymbol u'}=A_{\boldsymbol u}$ (because there are $q^2-1$ choices for the first basis element, $q^2-q$ choices for the second basis element and $|A_{\boldsymbol u}|=q^2$ choices for the offset).
\end{itemize}
Altogether,
\[\mathrm{REJ}_{2,d}^{\mathrm{MRM}}(T)\geq\frac{\Omega(q^6)}{q^{3t}}\cdot\min\left\{\frac\delta{3q},\frac1{2(t+1)q^{t+1}}\right\}=\Omega\left(\min\left\{q^{-3t+5}\delta,\frac{q^{-4t+5}}t\right\}\right).\]
We notice that $t\leq\lceil q(s-1/p)/(q-q/p)\rceil\leq\lceil s/(1-1/p)\rceil\leq2s$. Thus $\alpha=\Omega(q^{-6s+5})$ and $c=\Omega(q^{-8s+4})$ as promised.\hfill$\triangleleft$

\begin{thebibliography}{20}
\bibitem[1]{ref1} Noga Alon. Combinatorial nullstellensatz. Combinatorics, Probability and Computing, 8(1- 2):7–29, 1999.
\bibitem[2]{ref2} Noga Alon, Tali Kaufman, Michael Krivelevich, Simon Litsyn, and Dana Ron. Testing reed-muller codes. IEEE Transactions on Information Theory, 51(11):4032–4039, 2005.
\bibitem[3]{ref3} Simeon Ball and Oriol Serra. Punctured combinatorial nullstellensätze. Combinatorica, 29(5):511–522, 2009.
\bibitem[4]{ref4} Arnab Bhattacharyya, Swastik Kopparty, Grant Schoenebeck, Madhu Sudan, and David Zuckerman. Optimal testing of reed-muller codes. In 2010 IEEE 51st Annual Symposium on Foundations of Computer Science, pages 488–497. IEEE, 2010.
\bibitem[5]{ref5} David Cox, John Little, and Donal OShea. Ideals, varieties, and algorithms: an introduction to computational algebraic geometry and commutative algebra. Springer Science \& Business Media, 2013.
\bibitem[6]{ref6} Zeev Dvir, Swastik Kopparty, Shubhangi Saraf, and Madhu Sudan. Extensions to the method of multiplicities, with applications to kakeya sets and mergers. SIAM Journal on Computing, 42(6):2305–2328, 2013.
\bibitem[7]{ref7} Katalin Friedl and Madhu Sudan. Some improvements to total degree tests. In Proceedings Third Israel Symposium on the Theory of Computing and Systems, pages 190–198. IEEE, 1995.
\bibitem[8]{ref8} Peter Gemmell, Richard J. Lipton, Ronitt Rubinfeld, Madhu Sudan, and Avi Wigderson. Self-testing/correcting for polynomials and for approximate functions. In Cris Koutsougeras and Jeffrey Scott Vitter, editors, Proceedings of the 23rd Annual ACM Symposium on Theory of Computing, May 5-8, 1991, New Orleans, Louisiana, USA, pages 32–42. ACM, 1991. doi:10.1145/103418.103429.
\bibitem[9]{ref9} Venkatesan Guruswami and Carol Wang. Linear-algebraic list decoding for variants of reed– solomon codes. IEEE Transactions on Information Theory, 59(6):3257–3268, 2013.
\bibitem[10]{ref10} Elad Haramaty, Amir Shpilka, and Madhu Sudan. Optimal testing of multivariate polynomials over small prime fields. SIAM Journal on Computing, 42(2):536–562, 2013.
\bibitem[11]{ref11} Charanjit S Jutla, Anindya C Patthak, Atri Rudra, and David Zuckerman. Testing low-degree polynomials over prime fields. In 45th Annual IEEE Symposium on Foundations of Computer Science, pages 423–432. IEEE, 2004.
\bibitem[12]{ref12} Tali Kaufman and Dana Ron. Testing polynomials over general fields. SIAM Journal on Computing, 36(3):779–802, 2006.
\bibitem[13]{ref13} Swastik Kopparty. Some remarks on multiplicity codes. Discrete Geometry and Algebraic Combinatorics, 625:155–176, 2013.
\bibitem[14]{ref14} Swastik Kopparty, Or Meir, Noga Ron-Zewi, and Shubhangi Saraf. High-rate locally correctable and locally testable codes with sub-polynomial query complexity. Journal of the ACM (JACM), 64(2):1–42, 2017.
\bibitem[15]{ref15} Swastik Kopparty, Shubhangi Saraf, and Sergey Yekhanin. High-rate codes with sublinear-time decoding. Journal of the ACM (JACM), 61(5):1–20, 2014.
\bibitem[16]{ref16} Rasmus Refslund Nielsen. List decoding of linear block codes. PhD thesis, DTU, 2001.
\bibitem[17]{ref17} M Yu Rosenbloom and Michael Anatol’evich Tsfasman. Codes for the m-metric. Problemy Peredachi Informatsii, 33(1):55–63, 1997.
\bibitem[18]{ref18} Ronitt Rubinfeld and Madhu Sudan. Self-testing polynomial functions efficiently and over rational domains. In Greg N. Frederickson, editor, Proceedings of the Third Annual ACM/SIGACT-SIAM Symposium on Discrete Algorithms, 27-29 January 1992, Orlando, Florida, USA, pages 23–32. ACM/SIAM, 1992. URL: http://dl.acm.org/citation.cfm?id=139404.139410.
\bibitem[19]{ref19} Ronitt Rubinfeld and Madhu Sudan. Robust characterizations of polynomials with applications to program testing. SIAM Journal on Computing, 25(2):252–271, 1996.
\bibitem[20]{ref20} Madhu Sudan, Luca Trevisan, and Salil Vadhan. Pseudorandom generators without the xor lemma. Journal of Computer and System Sciences, 62(2):236–266, 2001.
\end{thebibliography}
\end{document}
